% Section S5 (temporal super-resolution results). Every number is a \Res, \ResDiff, \ResAbsDiff or \ResRatio call.
\section{Temporal super-resolution: results}\label{sec:tsrr}

This section reports every output of the temporal super-resolution (TSR) notebook, \texttt{notebooks/neural-supersampling.py}, from its saved metric files, plots, predictions and logs. For each output it states what is measured, what it shows, the mechanism that the method (\cref{sec:tsrm}) suggests, and what it does not show. The notebook trained one set of networks in one run, from \Res{run.neural-supersampling.start} to \Res{run.neural-supersampling.end}. Every number here is generated from the notebook's files by \texttt{tools/build\_paper\_results.py}, and every difference is computed by the same script. Each claim carries a status label as defined in \cref{sec:status-legend}. Because the saved outputs come from one run, a result that depends on training or evaluation noise is tagged \Preliminary. Only deterministic results, such as cost counts, are tagged \Confirmed.

\subsection{Protocol and reading rules}\label{sec:tsrr-protocol}

\emph{Test suites.} The test suites are built from the test split only. Every suite is evaluated at 3x, and the 2x student is evaluated on the suites of \cref{sec:tsrr-2x}.
\begin{description}
\item[Synthetic-game] Layered scenes rendered from held-out game stills and DIV2K validation images \cite{agustsson2017div2k} at $384\times384$ output, with six stills per motion class. The four classes are static, slow (0.3 to 3 output pixels per frame), medium (3 to 10) and fast (10 to 28), as defined in the notebook. Motion, disocclusion, thin geometry and particles are exact.
\item[TartanAir] Two held-out environments with exact optical flow and depth, at $576\times480$ \cite{wang2020tartanair}.
\item[Sintel] The \texttt{market} and \texttt{cave} scene families, with exact flow, at $1008\times432$ \cite{butler2012sintel}.
\item[Vid4 and Vimeo-1080p] Real clips with RAFT motion \cite{teed2020raft}, computed from the clean frames. Vid4 \cite{liu2014vid4} has four clips. The Vimeo-1080p clips are held-out test clips of 1080p Vimeo videos\footnote{Hugging Face dataset \texttt{danjacobellis/vimeo1080p}, \url{https://huggingface.co/datasets/danjacobellis/vimeo1080p}, accessed 9 November 2026. No licence is stated.} taken from the dataset's validation shard, which the pipeline assigns to the test split. Two blank clips were skipped.
\item[GameIR native] Held-out town 05 of the GameIR dataset \cite{zhou2024gameir}, whose clips the notebook describes as CARLA and Unreal Engine 4 renders \cite{dosovitskiy2017carla}. The 720p native render is the input and the 1440p render the target, so this is the only suite whose input is a real renderer's low-resolution image. Its clips have six frames, so the metrics use five frames per clip, and LPIPS and tOF use one frame per clip (frame 4).
\end{description}

\emph{Modes.} In \emph{engine} mode the input is a jittered single-sample render with motion vectors, and with depth where the suite has it. On the synthetic, TartanAir and Sintel suites the motion vectors are the exact motion of the scene. On Vid4, Vimeo-1080p and GameIR they come from RAFT on the clean full-resolution frames. This makes the motion of the video suites an oracle in engine mode, and the analytic baseline gets the same motion. In \emph{plain} mode the input is a bicubic downscale with no jitter, depth or motion vectors, and the networks use their built-in motion estimator. The analytic TAA baseline needs motion vectors, so it is not run in plain mode.

\emph{Metrics.} The definitions of the image and temporal metrics are in \cref{sec:bg-metrics}. The code computes the following for every frame. The tables average them over frames 1 to 15 of each 16-frame clip, so the first frame, which has no history, is skipped.
\begin{description}
\item[PSNR-Y and SSIM-Y] Peak signal-to-noise ratio and structural similarity \cite{wang2004ssim} on luma, with a border of one scale factor cropped on every side.
\item[LPIPS] The learned perceptual distance \cite{zhang2018lpips}, computed on every fourth frame (frames 4, 8 and 12 of a 16-frame clip).
\item[Disocc., Detail and Part.] PSNR-Y inside the disoccluded pixels (where the history is invalid), inside the top 15 per cent of ground-truth gradient magnitude (thin geometry, foliage and text), and inside the particle mask. The particle mask exists only for the synthetic scenes, so the real suites show n/a.
\item[Temp.] The temporal-change error. It is the mean absolute difference between the change of the output and the change of the ground truth between consecutive frames, after warping by the ground-truth motion, over the pixels visible in both frames.
\item[Warp] The warping error. It is the mean absolute difference between each output and the previous output warped by the ground-truth motion. It is computed only where the motion is exact, so the RAFT-motion suites show n/a.
\item[tOF] The temporal optical-flow error \cite{chu2020tecogan}: the mean length, in pixels, of the difference between the RAFT flow of the output and the RAFT flow of the ground truth, on every fourth frame.
\end{description}
Higher is better for PSNR and SSIM, and lower is better for LPIPS, Temp., Warp and tOF.

Two properties of the metrics limit their reading. Temp. and Warp reward blur, because a smooth output changes little from frame to frame, so a stability metric is read together with PSNR-Y. The real suites average over four or eight clips, so differences of a few hundredths of a decibel, or a few thousandths of LPIPS, are not resolved.

\subsection{Training runs}\label{sec:tsrr-training}

\emph{Budgets and stage times.} \Cref{tab:tsrr-stages} compares the planned budgets in the configuration with the measured stage times. The planned budgets were not scaled down, so every stage received its planned time. The \tier{igpu} 3x and \tier{low} 2x students shared one stage, and the three ablation pairs ran one after another. \Confirmed{} The run stayed inside its \Res{tsr.config.budget_h}-hour ceiling. The notebook logs the total run time as \Res{run.neural-supersampling.hours} hours at the end of evaluation, and its summary cell records \Res{tsr.summary.run_hours} hours, because the summary cell runs after the ONNX export.

\begin{table}[!htbp]
\caption{Planned and measured stage times of the TSR run, in minutes. The planned values come from the notebook configuration and the measured values from its stage timer. The ablations are three pairs of runs with a 30-minute budget each, and the caches stage had a 30-minute decoding budget}\label{tab:tsrr-stages}
\begin{tabular*}{\columnwidth}{@{\extracolsep\fill}lrr}
\toprule
Stage & Planned & Measured \\
\midrule
Teacher, both GPUs & \Res{tsr.config.cfg.minutes.teacher} & \Res{tsr.stage_times.train_teacher.minutes} \\
\tier{low} 3x student, both GPUs & \Res{tsr.config.cfg.minutes.student} & \Res{tsr.stage_times.train_student.minutes} \\
\tier{igpu} 3x and \tier{low} 2x, in parallel & \Res{tsr.config.cfg.minutes.pair} & \Res{tsr.stage_times.train_pair.minutes} \\
Ablations, three pairs & \Res{tsr.config.cfg.minutes.ablation} each & \Res{tsr.stage_times.ablations.minutes} \\
Training caches & \Res{tsr.config.cfg.cache.build_min} (budget) & \Res{tsr.stage_times.caches.minutes} \\
Batch probe & n/a & \Res{tsr.stage_times.batch_probe.minutes} \\
Test suites and evaluation & n/a & \Res{tsr.stage_times.test_suites.minutes} and \Res{tsr.stage_times.evaluation.minutes} \\
\botrule
\end{tabular*}
\end{table}

\emph{Hardware and batch.} The run had \Res{tsr.hardware.gpu_list.0.name} GPUs with \Res{tsr.hardware.gpu_list.0.vram_gb} GB each, \Res{tsr.hardware.cpu_count} CPU cores, \Res{tsr.hardware.ram_total_gb} GB of RAM, PyTorch \Res{tsr.hardware.torch} and CUDA \Res{tsr.hardware.cuda}. The batch probe ran six-frame training steps with a Charbonnier loss only, on GPU 0, for batch sizes from 4 to 48 sequences. It therefore omits the gradient, Fourier, temporal, motion and perceptual terms of the real loss, and the notebook counts this as the margin. Every tested size fitted. The teacher used \Res{tsr.batch_probe.teacher.48.peak_vram_pct} per cent of VRAM at batch 48, and the student \Res{tsr.batch_probe.student.48.peak_vram_pct} per cent. The teacher's throughput levelled off early, at \Res{tsr.batch_probe.teacher.32.seq_per_s} sequences per second at batch 32 and \Res{tsr.batch_probe.teacher.48.seq_per_s} at batch 48. The chosen batch of 48 is the largest size in the list, not a memory limit, because the 75 per cent VRAM rule never bound. \Observation{} A larger batch would probably have fitted, but the probe did not test one.

\emph{Validation and failed runs.} \Cref{tab:tsrr-bestval} gives the best validation round of every run. The checkpoint metric is the mean of the render and plain-video validation scores, which is the mean of the two middle columns. The validation batches are built with each run's own recipe, and the tiers differ in scale, so the validation values are comparable only within a group of the same scale and recipe. No run failed, and the failed-run record is empty.

\begin{table}[!htbp]
\caption{Best validation round of every training run. PSNR-Y is the checkpoint metric (mean of the render and video copies). Last frame is the PSNR-Y of the last validation frame at the same round. Temp. err. is the validation temporal-change error. Samples/s and Peak GB are the mean throughput and the peak memory of the run. Values are comparable only within a scale and recipe}\label{tab:tsrr-bestval}
\footnotesize
\setlength{\tabcolsep}{2.5pt}
\begin{tabular*}{\textwidth}{@{\extracolsep\fill}lrrrrrrrr}
\toprule
Run & Round & PSNR-Y & Render & Video & Last & Temp. & Samples/s & GB \\
\midrule
Teacher 3x & \Res{tsr.best_validation.tsr_teacher_s3.best_round} & \Res{tsr.best_validation.tsr_teacher_s3.val_psnr_y} & \Res{tsr.best_validation.tsr_teacher_s3.val_psnr_y_render} & \Res{tsr.best_validation.tsr_teacher_s3.val_psnr_y_video} & \Res{tsr.best_validation.tsr_teacher_s3.val_psnr_y_last} & \Res{tsr.best_validation.tsr_teacher_s3.val_tce} & \Res{tsr.best_validation.tsr_teacher_s3.samples_per_s} & \Res{tsr.best_validation.tsr_teacher_s3.peak_mem_gb} \\
\tier{low} 3x & \Res{tsr.best_validation.tsr_low_s3.best_round} & \Res{tsr.best_validation.tsr_low_s3.val_psnr_y} & \Res{tsr.best_validation.tsr_low_s3.val_psnr_y_render} & \Res{tsr.best_validation.tsr_low_s3.val_psnr_y_video} & \Res{tsr.best_validation.tsr_low_s3.val_psnr_y_last} & \Res{tsr.best_validation.tsr_low_s3.val_tce} & \Res{tsr.best_validation.tsr_low_s3.samples_per_s} & \Res{tsr.best_validation.tsr_low_s3.peak_mem_gb} \\
\tier{igpu} 3x & \Res{tsr.best_validation.tsr_igpu_s3.best_round} & \Res{tsr.best_validation.tsr_igpu_s3.val_psnr_y} & \Res{tsr.best_validation.tsr_igpu_s3.val_psnr_y_render} & \Res{tsr.best_validation.tsr_igpu_s3.val_psnr_y_video} & \Res{tsr.best_validation.tsr_igpu_s3.val_psnr_y_last} & \Res{tsr.best_validation.tsr_igpu_s3.val_tce} & \Res{tsr.best_validation.tsr_igpu_s3.samples_per_s} & \Res{tsr.best_validation.tsr_igpu_s3.peak_mem_gb} \\
\tier{low} 2x & \Res{tsr.best_validation.tsr_low_s2.best_round} & \Res{tsr.best_validation.tsr_low_s2.val_psnr_y} & \Res{tsr.best_validation.tsr_low_s2.val_psnr_y_render} & \Res{tsr.best_validation.tsr_low_s2.val_psnr_y_video} & \Res{tsr.best_validation.tsr_low_s2.val_psnr_y_last} & \Res{tsr.best_validation.tsr_low_s2.val_tce} & \Res{tsr.best_validation.tsr_low_s2.samples_per_s} & \Res{tsr.best_validation.tsr_low_s2.peak_mem_gb} \\
Reference & \Res{tsr.best_validation.abl_reference.best_round} & \Res{tsr.best_validation.abl_reference.val_psnr_y} & \Res{tsr.best_validation.abl_reference.val_psnr_y_render} & \Res{tsr.best_validation.abl_reference.val_psnr_y_video} & \Res{tsr.best_validation.abl_reference.val_psnr_y_last} & \Res{tsr.best_validation.abl_reference.val_tce} & \Res{tsr.best_validation.abl_reference.samples_per_s} & \Res{tsr.best_validation.abl_reference.peak_mem_gb} \\
No history & \Res{tsr.best_validation.abl_no_history.best_round} & \Res{tsr.best_validation.abl_no_history.val_psnr_y} & \Res{tsr.best_validation.abl_no_history.val_psnr_y_render} & \Res{tsr.best_validation.abl_no_history.val_psnr_y_video} & \Res{tsr.best_validation.abl_no_history.val_psnr_y_last} & \Res{tsr.best_validation.abl_no_history.val_tce} & \Res{tsr.best_validation.abl_no_history.samples_per_s} & \Res{tsr.best_validation.abl_no_history.peak_mem_gb} \\
No G-buf. & \Res{tsr.best_validation.abl_no_gbuffers.best_round} & \Res{tsr.best_validation.abl_no_gbuffers.val_psnr_y} & \Res{tsr.best_validation.abl_no_gbuffers.val_psnr_y_render} & \Res{tsr.best_validation.abl_no_gbuffers.val_psnr_y_video} & \Res{tsr.best_validation.abl_no_gbuffers.val_psnr_y_last} & \Res{tsr.best_validation.abl_no_gbuffers.val_tce} & \Res{tsr.best_validation.abl_no_gbuffers.samples_per_s} & \Res{tsr.best_validation.abl_no_gbuffers.peak_mem_gb} \\
No motion & \Res{tsr.best_validation.abl_no_motion_net.best_round} & \Res{tsr.best_validation.abl_no_motion_net.val_psnr_y} & \Res{tsr.best_validation.abl_no_motion_net.val_psnr_y_render} & \Res{tsr.best_validation.abl_no_motion_net.val_psnr_y_video} & \Res{tsr.best_validation.abl_no_motion_net.val_psnr_y_last} & \Res{tsr.best_validation.abl_no_motion_net.val_tce} & \Res{tsr.best_validation.abl_no_motion_net.samples_per_s} & \Res{tsr.best_validation.abl_no_motion_net.peak_mem_gb} \\
No video & \Res{tsr.best_validation.abl_no_video_training.best_round} & \Res{tsr.best_validation.abl_no_video_training.val_psnr_y} & \Res{tsr.best_validation.abl_no_video_training.val_psnr_y_render} & \Res{tsr.best_validation.abl_no_video_training.val_psnr_y_video} & \Res{tsr.best_validation.abl_no_video_training.val_psnr_y_last} & \Res{tsr.best_validation.abl_no_video_training.val_tce} & \Res{tsr.best_validation.abl_no_video_training.samples_per_s} & \Res{tsr.best_validation.abl_no_video_training.peak_mem_gb} \\
No unroll & \Res{tsr.best_validation.abl_no_long_unroll.best_round} & \Res{tsr.best_validation.abl_no_long_unroll.val_psnr_y} & \Res{tsr.best_validation.abl_no_long_unroll.val_psnr_y_render} & \Res{tsr.best_validation.abl_no_long_unroll.val_psnr_y_video} & \Res{tsr.best_validation.abl_no_long_unroll.val_psnr_y_last} & \Res{tsr.best_validation.abl_no_long_unroll.val_tce} & \Res{tsr.best_validation.abl_no_long_unroll.samples_per_s} & \Res{tsr.best_validation.abl_no_long_unroll.peak_mem_gb} \\
\botrule
\end{tabular*}
\end{table}

\Cref{fig:tsrr-training} shows the validation curves. The validation PSNR-Y of the teacher and of the \tier{low} 3x student keeps rising over the budget, and the curves flatten at the end, where the cosine schedule ends. The 2x student has the highest validation score, as expected for the easier 2x task, so the scales are not compared directly. The temporal-change error of the students falls and then levels off, and the teacher's training loss ends lowest. \Observation{} The schedules ended while the models were still improving, so these numbers measure the budget as much as the architecture. The 30-minute ablations stop earlier, so the ablation results describe short-budget training.

\begin{figure}[!htbp]
\centering
\includegraphics[width=\textwidth]{\TsrOut/plots/training_curves.png}
\caption{Validation PSNR-Y (left), validation temporal-change error (middle) and training loss (right) against wall-clock minutes for the teacher, the main students, the 2x student and the six 30-minute ablation runs (dashed in the first two panels). The training loss is shown for the main stages only}\label{fig:tsrr-training}
\end{figure}

\subsection{Main results at 3x}\label{sec:tsrr-main}

\Cref{tab:tsrr-main-synthetic-game,tab:tsrr-main-tartanair} give every metric at 3x for each suite and mode. The rows are the analytic TAA baseline (engine mode only), the bicubic and bilinear upsamplers, the trained tiers and the teacher. The ablations are in \cref{sec:tsrr-ablation}, and \cref{sec:tsrr-main-summary} compares the suites.

\begin{table}[!htbp]
\caption{Synthetic-game suite at 3x. Means over per-frame rows (frames 1 to 15 of each clip). Engine mode: \Res{tsr.test_per_frame.synthetic-game.engine.3.student_s3.n} rows per method. Plain mode: \Res{tsr.test_per_frame.synthetic-game.plain.3.student_s3.n} rows per method. Disocc., Detail and Part.: PSNR-Y inside disoccluded, fine-detail and particle pixels. Temp.: temporal-change error. Warp: warping error. tOF: temporal optical-flow error in pixels. n/a: not computed for this suite.}\label{tab:tsrr-main-synthetic-game}
\footnotesize
\setlength{\tabcolsep}{2.5pt}
\begin{tabular*}{\textwidth}{@{\extracolsep\fill}lrrrrrrrrr}
\toprule
Method & PSNR-Y & SSIM-Y & LPIPS & Disocc. & Detail & Part. & Temp. & tOF & Warp \\
\midrule
\multicolumn{10}{@{}l}{\emph{Engine mode}} \\
Bicubic & \Res{tsr.test_results.synthetic-game.engine.3.bicubic.psnr_y} & \Res{tsr.test_results.synthetic-game.engine.3.bicubic.ssim_y} & \Res{tsr.test_results.synthetic-game.engine.3.bicubic.lpips} & \Res{tsr.test_results.synthetic-game.engine.3.bicubic.psnr_disocc} & \Res{tsr.test_results.synthetic-game.engine.3.bicubic.psnr_detail} & \Res{tsr.test_results.synthetic-game.engine.3.bicubic.psnr_particles} & \Res{tsr.test_results.synthetic-game.engine.3.bicubic.temporal_err} & \Res{tsr.test_results.synthetic-game.engine.3.bicubic.tof} & \Res{tsr.test_results.synthetic-game.engine.3.bicubic.warp_err} \\
Bilinear & \Res{tsr.test_results.synthetic-game.engine.3.bilinear.psnr_y} & \Res{tsr.test_results.synthetic-game.engine.3.bilinear.ssim_y} & \Res{tsr.test_results.synthetic-game.engine.3.bilinear.lpips} & \Res{tsr.test_results.synthetic-game.engine.3.bilinear.psnr_disocc} & \Res{tsr.test_results.synthetic-game.engine.3.bilinear.psnr_detail} & \Res{tsr.test_results.synthetic-game.engine.3.bilinear.psnr_particles} & \Res{tsr.test_results.synthetic-game.engine.3.bilinear.temporal_err} & \Res{tsr.test_results.synthetic-game.engine.3.bilinear.tof} & \Res{tsr.test_results.synthetic-game.engine.3.bilinear.warp_err} \\
Analytic TAA & \Res{tsr.test_results.synthetic-game.engine.3.analytic_taau.psnr_y} & \Res{tsr.test_results.synthetic-game.engine.3.analytic_taau.ssim_y} & \Res{tsr.test_results.synthetic-game.engine.3.analytic_taau.lpips} & \Res{tsr.test_results.synthetic-game.engine.3.analytic_taau.psnr_disocc} & \Res{tsr.test_results.synthetic-game.engine.3.analytic_taau.psnr_detail} & \Res{tsr.test_results.synthetic-game.engine.3.analytic_taau.psnr_particles} & \Res{tsr.test_results.synthetic-game.engine.3.analytic_taau.temporal_err} & \Res{tsr.test_results.synthetic-game.engine.3.analytic_taau.tof} & \Res{tsr.test_results.synthetic-game.engine.3.analytic_taau.warp_err} \\
\tier{igpu} & \Res{tsr.test_results.synthetic-game.engine.3.igpu_s3.psnr_y} & \Res{tsr.test_results.synthetic-game.engine.3.igpu_s3.ssim_y} & \Res{tsr.test_results.synthetic-game.engine.3.igpu_s3.lpips} & \Res{tsr.test_results.synthetic-game.engine.3.igpu_s3.psnr_disocc} & \Res{tsr.test_results.synthetic-game.engine.3.igpu_s3.psnr_detail} & \Res{tsr.test_results.synthetic-game.engine.3.igpu_s3.psnr_particles} & \Res{tsr.test_results.synthetic-game.engine.3.igpu_s3.temporal_err} & \Res{tsr.test_results.synthetic-game.engine.3.igpu_s3.tof} & \Res{tsr.test_results.synthetic-game.engine.3.igpu_s3.warp_err} \\
\tier{low} student & \Res{tsr.test_results.synthetic-game.engine.3.student_s3.psnr_y} & \Res{tsr.test_results.synthetic-game.engine.3.student_s3.ssim_y} & \Res{tsr.test_results.synthetic-game.engine.3.student_s3.lpips} & \Res{tsr.test_results.synthetic-game.engine.3.student_s3.psnr_disocc} & \Res{tsr.test_results.synthetic-game.engine.3.student_s3.psnr_detail} & \Res{tsr.test_results.synthetic-game.engine.3.student_s3.psnr_particles} & \Res{tsr.test_results.synthetic-game.engine.3.student_s3.temporal_err} & \Res{tsr.test_results.synthetic-game.engine.3.student_s3.tof} & \Res{tsr.test_results.synthetic-game.engine.3.student_s3.warp_err} \\
Teacher & \Res{tsr.test_results.synthetic-game.engine.3.teacher_s3.psnr_y} & \Res{tsr.test_results.synthetic-game.engine.3.teacher_s3.ssim_y} & \Res{tsr.test_results.synthetic-game.engine.3.teacher_s3.lpips} & \Res{tsr.test_results.synthetic-game.engine.3.teacher_s3.psnr_disocc} & \Res{tsr.test_results.synthetic-game.engine.3.teacher_s3.psnr_detail} & \Res{tsr.test_results.synthetic-game.engine.3.teacher_s3.psnr_particles} & \Res{tsr.test_results.synthetic-game.engine.3.teacher_s3.temporal_err} & \Res{tsr.test_results.synthetic-game.engine.3.teacher_s3.tof} & \Res{tsr.test_results.synthetic-game.engine.3.teacher_s3.warp_err} \\
\midrule
\multicolumn{10}{@{}l}{\emph{Plain mode}} \\
Bicubic & \Res{tsr.test_results.synthetic-game.plain.3.bicubic.psnr_y} & \Res{tsr.test_results.synthetic-game.plain.3.bicubic.ssim_y} & \Res{tsr.test_results.synthetic-game.plain.3.bicubic.lpips} & \Res{tsr.test_results.synthetic-game.plain.3.bicubic.psnr_disocc} & \Res{tsr.test_results.synthetic-game.plain.3.bicubic.psnr_detail} & n/a & \Res{tsr.test_results.synthetic-game.plain.3.bicubic.temporal_err} & \Res{tsr.test_results.synthetic-game.plain.3.bicubic.tof} & \Res{tsr.test_results.synthetic-game.plain.3.bicubic.warp_err} \\
Bilinear & \Res{tsr.test_results.synthetic-game.plain.3.bilinear.psnr_y} & \Res{tsr.test_results.synthetic-game.plain.3.bilinear.ssim_y} & \Res{tsr.test_results.synthetic-game.plain.3.bilinear.lpips} & \Res{tsr.test_results.synthetic-game.plain.3.bilinear.psnr_disocc} & \Res{tsr.test_results.synthetic-game.plain.3.bilinear.psnr_detail} & n/a & \Res{tsr.test_results.synthetic-game.plain.3.bilinear.temporal_err} & \Res{tsr.test_results.synthetic-game.plain.3.bilinear.tof} & \Res{tsr.test_results.synthetic-game.plain.3.bilinear.warp_err} \\
\tier{igpu} & \Res{tsr.test_results.synthetic-game.plain.3.igpu_s3.psnr_y} & \Res{tsr.test_results.synthetic-game.plain.3.igpu_s3.ssim_y} & \Res{tsr.test_results.synthetic-game.plain.3.igpu_s3.lpips} & \Res{tsr.test_results.synthetic-game.plain.3.igpu_s3.psnr_disocc} & \Res{tsr.test_results.synthetic-game.plain.3.igpu_s3.psnr_detail} & n/a & \Res{tsr.test_results.synthetic-game.plain.3.igpu_s3.temporal_err} & \Res{tsr.test_results.synthetic-game.plain.3.igpu_s3.tof} & \Res{tsr.test_results.synthetic-game.plain.3.igpu_s3.warp_err} \\
\tier{low} student & \Res{tsr.test_results.synthetic-game.plain.3.student_s3.psnr_y} & \Res{tsr.test_results.synthetic-game.plain.3.student_s3.ssim_y} & \Res{tsr.test_results.synthetic-game.plain.3.student_s3.lpips} & \Res{tsr.test_results.synthetic-game.plain.3.student_s3.psnr_disocc} & \Res{tsr.test_results.synthetic-game.plain.3.student_s3.psnr_detail} & n/a & \Res{tsr.test_results.synthetic-game.plain.3.student_s3.temporal_err} & \Res{tsr.test_results.synthetic-game.plain.3.student_s3.tof} & \Res{tsr.test_results.synthetic-game.plain.3.student_s3.warp_err} \\
Teacher & \Res{tsr.test_results.synthetic-game.plain.3.teacher_s3.psnr_y} & \Res{tsr.test_results.synthetic-game.plain.3.teacher_s3.ssim_y} & \Res{tsr.test_results.synthetic-game.plain.3.teacher_s3.lpips} & \Res{tsr.test_results.synthetic-game.plain.3.teacher_s3.psnr_disocc} & \Res{tsr.test_results.synthetic-game.plain.3.teacher_s3.psnr_detail} & n/a & \Res{tsr.test_results.synthetic-game.plain.3.teacher_s3.temporal_err} & \Res{tsr.test_results.synthetic-game.plain.3.teacher_s3.tof} & \Res{tsr.test_results.synthetic-game.plain.3.teacher_s3.warp_err} \\
\botrule
\end{tabular*}
\end{table}
\begin{table}[!htbp]
\caption{TartanAir suite at 3x. Means over per-frame rows (frames 1 to 15 of each clip). Engine mode: \Res{tsr.test_per_frame.tartanair.engine.3.student_s3.n} rows per method. Plain mode: \Res{tsr.test_per_frame.tartanair.plain.3.student_s3.n} rows per method. Disocc., Detail and Part.: PSNR-Y inside disoccluded, fine-detail and particle pixels. Temp.: temporal-change error. Warp: warping error. tOF: temporal optical-flow error in pixels. n/a: not computed for this suite.}\label{tab:tsrr-main-tartanair}
\footnotesize
\setlength{\tabcolsep}{2.5pt}
\begin{tabular*}{\textwidth}{@{\extracolsep\fill}lrrrrrrrrr}
\toprule
Method & PSNR-Y & SSIM-Y & LPIPS & Disocc. & Detail & Part. & Temp. & tOF & Warp \\
\midrule
\multicolumn{10}{@{}l}{\emph{Engine mode}} \\
Bicubic & \Res{tsr.test_results.tartanair.engine.3.bicubic.psnr_y} & \Res{tsr.test_results.tartanair.engine.3.bicubic.ssim_y} & \Res{tsr.test_results.tartanair.engine.3.bicubic.lpips} & \Res{tsr.test_results.tartanair.engine.3.bicubic.psnr_disocc} & \Res{tsr.test_results.tartanair.engine.3.bicubic.psnr_detail} & n/a & \Res{tsr.test_results.tartanair.engine.3.bicubic.temporal_err} & \Res{tsr.test_results.tartanair.engine.3.bicubic.tof} & \Res{tsr.test_results.tartanair.engine.3.bicubic.warp_err} \\
Bilinear & \Res{tsr.test_results.tartanair.engine.3.bilinear.psnr_y} & \Res{tsr.test_results.tartanair.engine.3.bilinear.ssim_y} & \Res{tsr.test_results.tartanair.engine.3.bilinear.lpips} & \Res{tsr.test_results.tartanair.engine.3.bilinear.psnr_disocc} & \Res{tsr.test_results.tartanair.engine.3.bilinear.psnr_detail} & n/a & \Res{tsr.test_results.tartanair.engine.3.bilinear.temporal_err} & \Res{tsr.test_results.tartanair.engine.3.bilinear.tof} & \Res{tsr.test_results.tartanair.engine.3.bilinear.warp_err} \\
Analytic TAA & \Res{tsr.test_results.tartanair.engine.3.analytic_taau.psnr_y} & \Res{tsr.test_results.tartanair.engine.3.analytic_taau.ssim_y} & \Res{tsr.test_results.tartanair.engine.3.analytic_taau.lpips} & \Res{tsr.test_results.tartanair.engine.3.analytic_taau.psnr_disocc} & \Res{tsr.test_results.tartanair.engine.3.analytic_taau.psnr_detail} & n/a & \Res{tsr.test_results.tartanair.engine.3.analytic_taau.temporal_err} & \Res{tsr.test_results.tartanair.engine.3.analytic_taau.tof} & \Res{tsr.test_results.tartanair.engine.3.analytic_taau.warp_err} \\
\tier{igpu} & \Res{tsr.test_results.tartanair.engine.3.igpu_s3.psnr_y} & \Res{tsr.test_results.tartanair.engine.3.igpu_s3.ssim_y} & \Res{tsr.test_results.tartanair.engine.3.igpu_s3.lpips} & \Res{tsr.test_results.tartanair.engine.3.igpu_s3.psnr_disocc} & \Res{tsr.test_results.tartanair.engine.3.igpu_s3.psnr_detail} & n/a & \Res{tsr.test_results.tartanair.engine.3.igpu_s3.temporal_err} & \Res{tsr.test_results.tartanair.engine.3.igpu_s3.tof} & \Res{tsr.test_results.tartanair.engine.3.igpu_s3.warp_err} \\
\tier{low} student & \Res{tsr.test_results.tartanair.engine.3.student_s3.psnr_y} & \Res{tsr.test_results.tartanair.engine.3.student_s3.ssim_y} & \Res{tsr.test_results.tartanair.engine.3.student_s3.lpips} & \Res{tsr.test_results.tartanair.engine.3.student_s3.psnr_disocc} & \Res{tsr.test_results.tartanair.engine.3.student_s3.psnr_detail} & n/a & \Res{tsr.test_results.tartanair.engine.3.student_s3.temporal_err} & \Res{tsr.test_results.tartanair.engine.3.student_s3.tof} & \Res{tsr.test_results.tartanair.engine.3.student_s3.warp_err} \\
Teacher & \Res{tsr.test_results.tartanair.engine.3.teacher_s3.psnr_y} & \Res{tsr.test_results.tartanair.engine.3.teacher_s3.ssim_y} & \Res{tsr.test_results.tartanair.engine.3.teacher_s3.lpips} & \Res{tsr.test_results.tartanair.engine.3.teacher_s3.psnr_disocc} & \Res{tsr.test_results.tartanair.engine.3.teacher_s3.psnr_detail} & n/a & \Res{tsr.test_results.tartanair.engine.3.teacher_s3.temporal_err} & \Res{tsr.test_results.tartanair.engine.3.teacher_s3.tof} & \Res{tsr.test_results.tartanair.engine.3.teacher_s3.warp_err} \\
\midrule
\multicolumn{10}{@{}l}{\emph{Plain mode}} \\
Bicubic & \Res{tsr.test_results.tartanair.plain.3.bicubic.psnr_y} & \Res{tsr.test_results.tartanair.plain.3.bicubic.ssim_y} & \Res{tsr.test_results.tartanair.plain.3.bicubic.lpips} & \Res{tsr.test_results.tartanair.plain.3.bicubic.psnr_disocc} & \Res{tsr.test_results.tartanair.plain.3.bicubic.psnr_detail} & n/a & \Res{tsr.test_results.tartanair.plain.3.bicubic.temporal_err} & \Res{tsr.test_results.tartanair.plain.3.bicubic.tof} & \Res{tsr.test_results.tartanair.plain.3.bicubic.warp_err} \\
Bilinear & \Res{tsr.test_results.tartanair.plain.3.bilinear.psnr_y} & \Res{tsr.test_results.tartanair.plain.3.bilinear.ssim_y} & \Res{tsr.test_results.tartanair.plain.3.bilinear.lpips} & \Res{tsr.test_results.tartanair.plain.3.bilinear.psnr_disocc} & \Res{tsr.test_results.tartanair.plain.3.bilinear.psnr_detail} & n/a & \Res{tsr.test_results.tartanair.plain.3.bilinear.temporal_err} & \Res{tsr.test_results.tartanair.plain.3.bilinear.tof} & \Res{tsr.test_results.tartanair.plain.3.bilinear.warp_err} \\
\tier{igpu} & \Res{tsr.test_results.tartanair.plain.3.igpu_s3.psnr_y} & \Res{tsr.test_results.tartanair.plain.3.igpu_s3.ssim_y} & \Res{tsr.test_results.tartanair.plain.3.igpu_s3.lpips} & \Res{tsr.test_results.tartanair.plain.3.igpu_s3.psnr_disocc} & \Res{tsr.test_results.tartanair.plain.3.igpu_s3.psnr_detail} & n/a & \Res{tsr.test_results.tartanair.plain.3.igpu_s3.temporal_err} & \Res{tsr.test_results.tartanair.plain.3.igpu_s3.tof} & \Res{tsr.test_results.tartanair.plain.3.igpu_s3.warp_err} \\
\tier{low} student & \Res{tsr.test_results.tartanair.plain.3.student_s3.psnr_y} & \Res{tsr.test_results.tartanair.plain.3.student_s3.ssim_y} & \Res{tsr.test_results.tartanair.plain.3.student_s3.lpips} & \Res{tsr.test_results.tartanair.plain.3.student_s3.psnr_disocc} & \Res{tsr.test_results.tartanair.plain.3.student_s3.psnr_detail} & n/a & \Res{tsr.test_results.tartanair.plain.3.student_s3.temporal_err} & \Res{tsr.test_results.tartanair.plain.3.student_s3.tof} & \Res{tsr.test_results.tartanair.plain.3.student_s3.warp_err} \\
Teacher & \Res{tsr.test_results.tartanair.plain.3.teacher_s3.psnr_y} & \Res{tsr.test_results.tartanair.plain.3.teacher_s3.ssim_y} & \Res{tsr.test_results.tartanair.plain.3.teacher_s3.lpips} & \Res{tsr.test_results.tartanair.plain.3.teacher_s3.psnr_disocc} & \Res{tsr.test_results.tartanair.plain.3.teacher_s3.psnr_detail} & n/a & \Res{tsr.test_results.tartanair.plain.3.teacher_s3.temporal_err} & \Res{tsr.test_results.tartanair.plain.3.teacher_s3.tof} & \Res{tsr.test_results.tartanair.plain.3.teacher_s3.warp_err} \\
\botrule
\end{tabular*}
\end{table}

\emph{Synthetic-game.} In engine mode the \tier{low} student reaches \Res{tsr.test_results.synthetic-game.engine.3.student_s3.psnr_y} dB PSNR-Y, which is \ResAbsDiff{tsr.test_results.synthetic-game.engine.3.student_s3.psnr_y}{tsr.test_results.synthetic-game.engine.3.bicubic.psnr_y} dB above bicubic and \ResAbsDiff{tsr.test_results.synthetic-game.engine.3.student_s3.psnr_y}{tsr.test_results.synthetic-game.engine.3.analytic_taau.psnr_y} dB above the analytic TAA baseline. The gain is largest in the hard regions. Disocclusion PSNR is \ResAbsDiff{tsr.test_results.synthetic-game.engine.3.student_s3.psnr_disocc}{tsr.test_results.synthetic-game.engine.3.analytic_taau.psnr_disocc} dB above TAA, and particle PSNR is \ResAbsDiff{tsr.test_results.synthetic-game.engine.3.student_s3.psnr_particles}{tsr.test_results.synthetic-game.engine.3.analytic_taau.psnr_particles} dB above TAA. Fine-detail PSNR is \ResAbsDiff{tsr.test_results.synthetic-game.engine.3.student_s3.psnr_detail}{tsr.test_results.synthetic-game.engine.3.bicubic.psnr_detail} dB above bicubic. \Preliminary{} The LPIPS of the student is \Res{tsr.test_results.synthetic-game.engine.3.student_s3.lpips} against \Res{tsr.test_results.synthetic-game.engine.3.bicubic.lpips} for bicubic, so its output is also perceptually closer to the target. The temporal-change error of the student is tied with TAA (\Res{tsr.test_results.synthetic-game.engine.3.student_s3.temporal_err} against \Res{tsr.test_results.synthetic-game.engine.3.analytic_taau.temporal_err}) and far below bicubic (\Res{tsr.test_results.synthetic-game.engine.3.bicubic.temporal_err}). A plausible mechanism is that TAA keeps its output stable by blending a fixed share of history, which also blurs it, while bicubic has no history and reproduces the sub-pixel jitter of each frame as flicker. The learned history weight keeps detail and stability together. \Observation{} \Cref{fig:tsrr-stability} plots this combination for every method.

\begin{figure}[!htbp]
\centering
\includegraphics[width=\textwidth]{\TsrOut/plots/quality_vs_stability.png}
\caption{Quality against stability at 3x. Left: engine mode, PSNR-Y (dB) against temporal-change error, one point per method and suite. Middle: the same for plain mode. Right: plain-video PSNR-Y of each method per suite. The scatter panels include the ablations, and their labels overlap}\label{fig:tsrr-stability}
\end{figure}

In plain mode the student reaches \Res{tsr.test_results.synthetic-game.plain.3.student_s3.psnr_y} dB, \ResAbsDiff{tsr.test_results.synthetic-game.plain.3.student_s3.psnr_y}{tsr.test_results.synthetic-game.plain.3.bicubic.psnr_y} dB above bicubic, and its LPIPS is \Res{tsr.test_results.synthetic-game.plain.3.student_s3.lpips} against \Res{tsr.test_results.synthetic-game.plain.3.bicubic.lpips}. Its temporal-change error is slightly below bicubic's (\Res{tsr.test_results.synthetic-game.plain.3.student_s3.temporal_err} against \Res{tsr.test_results.synthetic-game.plain.3.bicubic.temporal_err}), but its warping error is slightly above it (\Res{tsr.test_results.synthetic-game.plain.3.student_s3.warp_err} against \Res{tsr.test_results.synthetic-game.plain.3.bicubic.warp_err}). \Observation{} Without jitter the plain-video student gains little stability, and its warping error shows no gain.

\emph{TartanAir.} The student reaches \Res{tsr.test_results.tartanair.engine.3.student_s3.psnr_y} dB in engine mode, \ResAbsDiff{tsr.test_results.tartanair.engine.3.student_s3.psnr_y}{tsr.test_results.tartanair.engine.3.analytic_taau.psnr_y} dB above TAA and \ResAbsDiff{tsr.test_results.tartanair.engine.3.student_s3.psnr_y}{tsr.test_results.tartanair.engine.3.bicubic.psnr_y} dB above bicubic. Its temporal-change error is \Res{tsr.test_results.tartanair.engine.3.student_s3.temporal_err} against \Res{tsr.test_results.tartanair.engine.3.analytic_taau.temporal_err} for TAA. In plain mode the student reaches \Res{tsr.test_results.tartanair.plain.3.student_s3.psnr_y} dB against \Res{tsr.test_results.tartanair.plain.3.bicubic.psnr_y} dB for bicubic, and its LPIPS is \Res{tsr.test_results.tartanair.plain.3.student_s3.lpips} against \Res{tsr.test_results.tartanair.plain.3.bicubic.lpips}. \Preliminary{} The gains are of the same order on TartanAir as on the synthetic scenes in plain mode, and the visual comparisons of \cref{sec:tsrr-visual} show where they occur on this suite.

\Cref{tab:tsrr-main-sintel} gives the Sintel suite in the same layout.

\begin{table}[!htbp]
\caption{Sintel suite at 3x. Means over per-frame rows (frames 1 to 15 of each clip). Engine mode: \Res{tsr.test_per_frame.sintel.engine.3.student_s3.n} rows per method. Plain mode: \Res{tsr.test_per_frame.sintel.plain.3.student_s3.n} rows per method. Disocc., Detail and Part.: PSNR-Y inside disoccluded, fine-detail and particle pixels. Temp.: temporal-change error. Warp: warping error. tOF: temporal optical-flow error in pixels. n/a: not computed for this suite.}\label{tab:tsrr-main-sintel}
\footnotesize
\setlength{\tabcolsep}{2.5pt}
\begin{tabular*}{\textwidth}{@{\extracolsep\fill}lrrrrrrrrr}
\toprule
Method & PSNR-Y & SSIM-Y & LPIPS & Disocc. & Detail & Part. & Temp. & tOF & Warp \\
\midrule
\multicolumn{10}{@{}l}{\emph{Engine mode}} \\
Bicubic & \Res{tsr.test_results.sintel.engine.3.bicubic.psnr_y} & \Res{tsr.test_results.sintel.engine.3.bicubic.ssim_y} & \Res{tsr.test_results.sintel.engine.3.bicubic.lpips} & \Res{tsr.test_results.sintel.engine.3.bicubic.psnr_disocc} & \Res{tsr.test_results.sintel.engine.3.bicubic.psnr_detail} & n/a & \Res{tsr.test_results.sintel.engine.3.bicubic.temporal_err} & \Res{tsr.test_results.sintel.engine.3.bicubic.tof} & \Res{tsr.test_results.sintel.engine.3.bicubic.warp_err} \\
Bilinear & \Res{tsr.test_results.sintel.engine.3.bilinear.psnr_y} & \Res{tsr.test_results.sintel.engine.3.bilinear.ssim_y} & \Res{tsr.test_results.sintel.engine.3.bilinear.lpips} & \Res{tsr.test_results.sintel.engine.3.bilinear.psnr_disocc} & \Res{tsr.test_results.sintel.engine.3.bilinear.psnr_detail} & n/a & \Res{tsr.test_results.sintel.engine.3.bilinear.temporal_err} & \Res{tsr.test_results.sintel.engine.3.bilinear.tof} & \Res{tsr.test_results.sintel.engine.3.bilinear.warp_err} \\
Analytic TAA & \Res{tsr.test_results.sintel.engine.3.analytic_taau.psnr_y} & \Res{tsr.test_results.sintel.engine.3.analytic_taau.ssim_y} & \Res{tsr.test_results.sintel.engine.3.analytic_taau.lpips} & \Res{tsr.test_results.sintel.engine.3.analytic_taau.psnr_disocc} & \Res{tsr.test_results.sintel.engine.3.analytic_taau.psnr_detail} & n/a & \Res{tsr.test_results.sintel.engine.3.analytic_taau.temporal_err} & \Res{tsr.test_results.sintel.engine.3.analytic_taau.tof} & \Res{tsr.test_results.sintel.engine.3.analytic_taau.warp_err} \\
\tier{igpu} & \Res{tsr.test_results.sintel.engine.3.igpu_s3.psnr_y} & \Res{tsr.test_results.sintel.engine.3.igpu_s3.ssim_y} & \Res{tsr.test_results.sintel.engine.3.igpu_s3.lpips} & \Res{tsr.test_results.sintel.engine.3.igpu_s3.psnr_disocc} & \Res{tsr.test_results.sintel.engine.3.igpu_s3.psnr_detail} & n/a & \Res{tsr.test_results.sintel.engine.3.igpu_s3.temporal_err} & \Res{tsr.test_results.sintel.engine.3.igpu_s3.tof} & \Res{tsr.test_results.sintel.engine.3.igpu_s3.warp_err} \\
\tier{low} student & \Res{tsr.test_results.sintel.engine.3.student_s3.psnr_y} & \Res{tsr.test_results.sintel.engine.3.student_s3.ssim_y} & \Res{tsr.test_results.sintel.engine.3.student_s3.lpips} & \Res{tsr.test_results.sintel.engine.3.student_s3.psnr_disocc} & \Res{tsr.test_results.sintel.engine.3.student_s3.psnr_detail} & n/a & \Res{tsr.test_results.sintel.engine.3.student_s3.temporal_err} & \Res{tsr.test_results.sintel.engine.3.student_s3.tof} & \Res{tsr.test_results.sintel.engine.3.student_s3.warp_err} \\
Teacher & \Res{tsr.test_results.sintel.engine.3.teacher_s3.psnr_y} & \Res{tsr.test_results.sintel.engine.3.teacher_s3.ssim_y} & \Res{tsr.test_results.sintel.engine.3.teacher_s3.lpips} & \Res{tsr.test_results.sintel.engine.3.teacher_s3.psnr_disocc} & \Res{tsr.test_results.sintel.engine.3.teacher_s3.psnr_detail} & n/a & \Res{tsr.test_results.sintel.engine.3.teacher_s3.temporal_err} & \Res{tsr.test_results.sintel.engine.3.teacher_s3.tof} & \Res{tsr.test_results.sintel.engine.3.teacher_s3.warp_err} \\
\midrule
\multicolumn{10}{@{}l}{\emph{Plain mode}} \\
Bicubic & \Res{tsr.test_results.sintel.plain.3.bicubic.psnr_y} & \Res{tsr.test_results.sintel.plain.3.bicubic.ssim_y} & \Res{tsr.test_results.sintel.plain.3.bicubic.lpips} & \Res{tsr.test_results.sintel.plain.3.bicubic.psnr_disocc} & \Res{tsr.test_results.sintel.plain.3.bicubic.psnr_detail} & n/a & \Res{tsr.test_results.sintel.plain.3.bicubic.temporal_err} & \Res{tsr.test_results.sintel.plain.3.bicubic.tof} & \Res{tsr.test_results.sintel.plain.3.bicubic.warp_err} \\
Bilinear & \Res{tsr.test_results.sintel.plain.3.bilinear.psnr_y} & \Res{tsr.test_results.sintel.plain.3.bilinear.ssim_y} & \Res{tsr.test_results.sintel.plain.3.bilinear.lpips} & \Res{tsr.test_results.sintel.plain.3.bilinear.psnr_disocc} & \Res{tsr.test_results.sintel.plain.3.bilinear.psnr_detail} & n/a & \Res{tsr.test_results.sintel.plain.3.bilinear.temporal_err} & \Res{tsr.test_results.sintel.plain.3.bilinear.tof} & \Res{tsr.test_results.sintel.plain.3.bilinear.warp_err} \\
\tier{igpu} & \Res{tsr.test_results.sintel.plain.3.igpu_s3.psnr_y} & \Res{tsr.test_results.sintel.plain.3.igpu_s3.ssim_y} & \Res{tsr.test_results.sintel.plain.3.igpu_s3.lpips} & \Res{tsr.test_results.sintel.plain.3.igpu_s3.psnr_disocc} & \Res{tsr.test_results.sintel.plain.3.igpu_s3.psnr_detail} & n/a & \Res{tsr.test_results.sintel.plain.3.igpu_s3.temporal_err} & \Res{tsr.test_results.sintel.plain.3.igpu_s3.tof} & \Res{tsr.test_results.sintel.plain.3.igpu_s3.warp_err} \\
\tier{low} student & \Res{tsr.test_results.sintel.plain.3.student_s3.psnr_y} & \Res{tsr.test_results.sintel.plain.3.student_s3.ssim_y} & \Res{tsr.test_results.sintel.plain.3.student_s3.lpips} & \Res{tsr.test_results.sintel.plain.3.student_s3.psnr_disocc} & \Res{tsr.test_results.sintel.plain.3.student_s3.psnr_detail} & n/a & \Res{tsr.test_results.sintel.plain.3.student_s3.temporal_err} & \Res{tsr.test_results.sintel.plain.3.student_s3.tof} & \Res{tsr.test_results.sintel.plain.3.student_s3.warp_err} \\
Teacher & \Res{tsr.test_results.sintel.plain.3.teacher_s3.psnr_y} & \Res{tsr.test_results.sintel.plain.3.teacher_s3.ssim_y} & \Res{tsr.test_results.sintel.plain.3.teacher_s3.lpips} & \Res{tsr.test_results.sintel.plain.3.teacher_s3.psnr_disocc} & \Res{tsr.test_results.sintel.plain.3.teacher_s3.psnr_detail} & n/a & \Res{tsr.test_results.sintel.plain.3.teacher_s3.temporal_err} & \Res{tsr.test_results.sintel.plain.3.teacher_s3.tof} & \Res{tsr.test_results.sintel.plain.3.teacher_s3.warp_err} \\
\botrule
\end{tabular*}
\end{table}
\Cref{tab:tsrr-main-vid4} gives the Vid4 suite in the same layout.

\begin{table}[!htbp]
\caption{Vid4 suite at 3x. Means over per-frame rows (frames 1 to 15 of each clip). Engine mode: \Res{tsr.test_per_frame.vid4.engine.3.student_s3.n} rows per method. Plain mode: \Res{tsr.test_per_frame.vid4.plain.3.student_s3.n} rows per method. Disocc., Detail and Part.: PSNR-Y inside disoccluded, fine-detail and particle pixels. Temp.: temporal-change error. Warp: warping error. tOF: temporal optical-flow error in pixels. n/a: not computed for this suite.}\label{tab:tsrr-main-vid4}
\footnotesize
\setlength{\tabcolsep}{2.5pt}
\begin{tabular*}{\textwidth}{@{\extracolsep\fill}lrrrrrrrrr}
\toprule
Method & PSNR-Y & SSIM-Y & LPIPS & Disocc. & Detail & Part. & Temp. & tOF & Warp \\
\midrule
\multicolumn{10}{@{}l}{\emph{Engine mode}} \\
Bicubic & \Res{tsr.test_results.vid4.engine.3.bicubic.psnr_y} & \Res{tsr.test_results.vid4.engine.3.bicubic.ssim_y} & \Res{tsr.test_results.vid4.engine.3.bicubic.lpips} & \Res{tsr.test_results.vid4.engine.3.bicubic.psnr_disocc} & \Res{tsr.test_results.vid4.engine.3.bicubic.psnr_detail} & n/a & \Res{tsr.test_results.vid4.engine.3.bicubic.temporal_err} & \Res{tsr.test_results.vid4.engine.3.bicubic.tof} & n/a \\
Bilinear & \Res{tsr.test_results.vid4.engine.3.bilinear.psnr_y} & \Res{tsr.test_results.vid4.engine.3.bilinear.ssim_y} & \Res{tsr.test_results.vid4.engine.3.bilinear.lpips} & \Res{tsr.test_results.vid4.engine.3.bilinear.psnr_disocc} & \Res{tsr.test_results.vid4.engine.3.bilinear.psnr_detail} & n/a & \Res{tsr.test_results.vid4.engine.3.bilinear.temporal_err} & \Res{tsr.test_results.vid4.engine.3.bilinear.tof} & n/a \\
Analytic TAA & \Res{tsr.test_results.vid4.engine.3.analytic_taau.psnr_y} & \Res{tsr.test_results.vid4.engine.3.analytic_taau.ssim_y} & \Res{tsr.test_results.vid4.engine.3.analytic_taau.lpips} & \Res{tsr.test_results.vid4.engine.3.analytic_taau.psnr_disocc} & \Res{tsr.test_results.vid4.engine.3.analytic_taau.psnr_detail} & n/a & \Res{tsr.test_results.vid4.engine.3.analytic_taau.temporal_err} & \Res{tsr.test_results.vid4.engine.3.analytic_taau.tof} & n/a \\
\tier{igpu} & \Res{tsr.test_results.vid4.engine.3.igpu_s3.psnr_y} & \Res{tsr.test_results.vid4.engine.3.igpu_s3.ssim_y} & \Res{tsr.test_results.vid4.engine.3.igpu_s3.lpips} & \Res{tsr.test_results.vid4.engine.3.igpu_s3.psnr_disocc} & \Res{tsr.test_results.vid4.engine.3.igpu_s3.psnr_detail} & n/a & \Res{tsr.test_results.vid4.engine.3.igpu_s3.temporal_err} & \Res{tsr.test_results.vid4.engine.3.igpu_s3.tof} & n/a \\
\tier{low} student & \Res{tsr.test_results.vid4.engine.3.student_s3.psnr_y} & \Res{tsr.test_results.vid4.engine.3.student_s3.ssim_y} & \Res{tsr.test_results.vid4.engine.3.student_s3.lpips} & \Res{tsr.test_results.vid4.engine.3.student_s3.psnr_disocc} & \Res{tsr.test_results.vid4.engine.3.student_s3.psnr_detail} & n/a & \Res{tsr.test_results.vid4.engine.3.student_s3.temporal_err} & \Res{tsr.test_results.vid4.engine.3.student_s3.tof} & n/a \\
Teacher & \Res{tsr.test_results.vid4.engine.3.teacher_s3.psnr_y} & \Res{tsr.test_results.vid4.engine.3.teacher_s3.ssim_y} & \Res{tsr.test_results.vid4.engine.3.teacher_s3.lpips} & \Res{tsr.test_results.vid4.engine.3.teacher_s3.psnr_disocc} & \Res{tsr.test_results.vid4.engine.3.teacher_s3.psnr_detail} & n/a & \Res{tsr.test_results.vid4.engine.3.teacher_s3.temporal_err} & \Res{tsr.test_results.vid4.engine.3.teacher_s3.tof} & n/a \\
\midrule
\multicolumn{10}{@{}l}{\emph{Plain mode}} \\
Bicubic & \Res{tsr.test_results.vid4.plain.3.bicubic.psnr_y} & \Res{tsr.test_results.vid4.plain.3.bicubic.ssim_y} & \Res{tsr.test_results.vid4.plain.3.bicubic.lpips} & \Res{tsr.test_results.vid4.plain.3.bicubic.psnr_disocc} & \Res{tsr.test_results.vid4.plain.3.bicubic.psnr_detail} & n/a & \Res{tsr.test_results.vid4.plain.3.bicubic.temporal_err} & \Res{tsr.test_results.vid4.plain.3.bicubic.tof} & n/a \\
Bilinear & \Res{tsr.test_results.vid4.plain.3.bilinear.psnr_y} & \Res{tsr.test_results.vid4.plain.3.bilinear.ssim_y} & \Res{tsr.test_results.vid4.plain.3.bilinear.lpips} & \Res{tsr.test_results.vid4.plain.3.bilinear.psnr_disocc} & \Res{tsr.test_results.vid4.plain.3.bilinear.psnr_detail} & n/a & \Res{tsr.test_results.vid4.plain.3.bilinear.temporal_err} & \Res{tsr.test_results.vid4.plain.3.bilinear.tof} & n/a \\
\tier{igpu} & \Res{tsr.test_results.vid4.plain.3.igpu_s3.psnr_y} & \Res{tsr.test_results.vid4.plain.3.igpu_s3.ssim_y} & \Res{tsr.test_results.vid4.plain.3.igpu_s3.lpips} & \Res{tsr.test_results.vid4.plain.3.igpu_s3.psnr_disocc} & \Res{tsr.test_results.vid4.plain.3.igpu_s3.psnr_detail} & n/a & \Res{tsr.test_results.vid4.plain.3.igpu_s3.temporal_err} & \Res{tsr.test_results.vid4.plain.3.igpu_s3.tof} & n/a \\
\tier{low} student & \Res{tsr.test_results.vid4.plain.3.student_s3.psnr_y} & \Res{tsr.test_results.vid4.plain.3.student_s3.ssim_y} & \Res{tsr.test_results.vid4.plain.3.student_s3.lpips} & \Res{tsr.test_results.vid4.plain.3.student_s3.psnr_disocc} & \Res{tsr.test_results.vid4.plain.3.student_s3.psnr_detail} & n/a & \Res{tsr.test_results.vid4.plain.3.student_s3.temporal_err} & \Res{tsr.test_results.vid4.plain.3.student_s3.tof} & n/a \\
Teacher & \Res{tsr.test_results.vid4.plain.3.teacher_s3.psnr_y} & \Res{tsr.test_results.vid4.plain.3.teacher_s3.ssim_y} & \Res{tsr.test_results.vid4.plain.3.teacher_s3.lpips} & \Res{tsr.test_results.vid4.plain.3.teacher_s3.psnr_disocc} & \Res{tsr.test_results.vid4.plain.3.teacher_s3.psnr_detail} & n/a & \Res{tsr.test_results.vid4.plain.3.teacher_s3.temporal_err} & \Res{tsr.test_results.vid4.plain.3.teacher_s3.tof} & n/a \\
\botrule
\end{tabular*}
\end{table}

\FloatBarrier
\Cref{tab:tsrr-main-vimeo1080p} gives the Vimeo-1080p suite in the same layout.

\begin{table}[!htbp]
\caption{Vimeo-1080p suite at 3x. Means over per-frame rows (frames 1 to 15 of each clip). Engine mode: \Res{tsr.test_per_frame.vimeo1080p.engine.3.student_s3.n} rows per method. Plain mode: \Res{tsr.test_per_frame.vimeo1080p.plain.3.student_s3.n} rows per method. Disocc., Detail and Part.: PSNR-Y inside disoccluded, fine-detail and particle pixels. Temp.: temporal-change error. Warp: warping error. tOF: temporal optical-flow error in pixels. n/a: not computed for this suite.}\label{tab:tsrr-main-vimeo1080p}
\footnotesize
\setlength{\tabcolsep}{2.5pt}
\begin{tabular*}{\textwidth}{@{\extracolsep\fill}lrrrrrrrrr}
\toprule
Method & PSNR-Y & SSIM-Y & LPIPS & Disocc. & Detail & Part. & Temp. & tOF & Warp \\
\midrule
\multicolumn{10}{@{}l}{\emph{Engine mode}} \\
Bicubic & \Res{tsr.test_results.vimeo1080p.engine.3.bicubic.psnr_y} & \Res{tsr.test_results.vimeo1080p.engine.3.bicubic.ssim_y} & \Res{tsr.test_results.vimeo1080p.engine.3.bicubic.lpips} & \Res{tsr.test_results.vimeo1080p.engine.3.bicubic.psnr_disocc} & \Res{tsr.test_results.vimeo1080p.engine.3.bicubic.psnr_detail} & n/a & \Res{tsr.test_results.vimeo1080p.engine.3.bicubic.temporal_err} & \Res{tsr.test_results.vimeo1080p.engine.3.bicubic.tof} & n/a \\
Bilinear & \Res{tsr.test_results.vimeo1080p.engine.3.bilinear.psnr_y} & \Res{tsr.test_results.vimeo1080p.engine.3.bilinear.ssim_y} & \Res{tsr.test_results.vimeo1080p.engine.3.bilinear.lpips} & \Res{tsr.test_results.vimeo1080p.engine.3.bilinear.psnr_disocc} & \Res{tsr.test_results.vimeo1080p.engine.3.bilinear.psnr_detail} & n/a & \Res{tsr.test_results.vimeo1080p.engine.3.bilinear.temporal_err} & \Res{tsr.test_results.vimeo1080p.engine.3.bilinear.tof} & n/a \\
Analytic TAA & \Res{tsr.test_results.vimeo1080p.engine.3.analytic_taau.psnr_y} & \Res{tsr.test_results.vimeo1080p.engine.3.analytic_taau.ssim_y} & \Res{tsr.test_results.vimeo1080p.engine.3.analytic_taau.lpips} & \Res{tsr.test_results.vimeo1080p.engine.3.analytic_taau.psnr_disocc} & \Res{tsr.test_results.vimeo1080p.engine.3.analytic_taau.psnr_detail} & n/a & \Res{tsr.test_results.vimeo1080p.engine.3.analytic_taau.temporal_err} & \Res{tsr.test_results.vimeo1080p.engine.3.analytic_taau.tof} & n/a \\
\tier{igpu} & \Res{tsr.test_results.vimeo1080p.engine.3.igpu_s3.psnr_y} & \Res{tsr.test_results.vimeo1080p.engine.3.igpu_s3.ssim_y} & \Res{tsr.test_results.vimeo1080p.engine.3.igpu_s3.lpips} & \Res{tsr.test_results.vimeo1080p.engine.3.igpu_s3.psnr_disocc} & \Res{tsr.test_results.vimeo1080p.engine.3.igpu_s3.psnr_detail} & n/a & \Res{tsr.test_results.vimeo1080p.engine.3.igpu_s3.temporal_err} & \Res{tsr.test_results.vimeo1080p.engine.3.igpu_s3.tof} & n/a \\
\tier{low} student & \Res{tsr.test_results.vimeo1080p.engine.3.student_s3.psnr_y} & \Res{tsr.test_results.vimeo1080p.engine.3.student_s3.ssim_y} & \Res{tsr.test_results.vimeo1080p.engine.3.student_s3.lpips} & \Res{tsr.test_results.vimeo1080p.engine.3.student_s3.psnr_disocc} & \Res{tsr.test_results.vimeo1080p.engine.3.student_s3.psnr_detail} & n/a & \Res{tsr.test_results.vimeo1080p.engine.3.student_s3.temporal_err} & \Res{tsr.test_results.vimeo1080p.engine.3.student_s3.tof} & n/a \\
Teacher & \Res{tsr.test_results.vimeo1080p.engine.3.teacher_s3.psnr_y} & \Res{tsr.test_results.vimeo1080p.engine.3.teacher_s3.ssim_y} & \Res{tsr.test_results.vimeo1080p.engine.3.teacher_s3.lpips} & \Res{tsr.test_results.vimeo1080p.engine.3.teacher_s3.psnr_disocc} & \Res{tsr.test_results.vimeo1080p.engine.3.teacher_s3.psnr_detail} & n/a & \Res{tsr.test_results.vimeo1080p.engine.3.teacher_s3.temporal_err} & \Res{tsr.test_results.vimeo1080p.engine.3.teacher_s3.tof} & n/a \\
\midrule
\multicolumn{10}{@{}l}{\emph{Plain mode}} \\
Bicubic & \Res{tsr.test_results.vimeo1080p.plain.3.bicubic.psnr_y} & \Res{tsr.test_results.vimeo1080p.plain.3.bicubic.ssim_y} & \Res{tsr.test_results.vimeo1080p.plain.3.bicubic.lpips} & \Res{tsr.test_results.vimeo1080p.plain.3.bicubic.psnr_disocc} & \Res{tsr.test_results.vimeo1080p.plain.3.bicubic.psnr_detail} & n/a & \Res{tsr.test_results.vimeo1080p.plain.3.bicubic.temporal_err} & \Res{tsr.test_results.vimeo1080p.plain.3.bicubic.tof} & n/a \\
Bilinear & \Res{tsr.test_results.vimeo1080p.plain.3.bilinear.psnr_y} & \Res{tsr.test_results.vimeo1080p.plain.3.bilinear.ssim_y} & \Res{tsr.test_results.vimeo1080p.plain.3.bilinear.lpips} & \Res{tsr.test_results.vimeo1080p.plain.3.bilinear.psnr_disocc} & \Res{tsr.test_results.vimeo1080p.plain.3.bilinear.psnr_detail} & n/a & \Res{tsr.test_results.vimeo1080p.plain.3.bilinear.temporal_err} & \Res{tsr.test_results.vimeo1080p.plain.3.bilinear.tof} & n/a \\
\tier{igpu} & \Res{tsr.test_results.vimeo1080p.plain.3.igpu_s3.psnr_y} & \Res{tsr.test_results.vimeo1080p.plain.3.igpu_s3.ssim_y} & \Res{tsr.test_results.vimeo1080p.plain.3.igpu_s3.lpips} & \Res{tsr.test_results.vimeo1080p.plain.3.igpu_s3.psnr_disocc} & \Res{tsr.test_results.vimeo1080p.plain.3.igpu_s3.psnr_detail} & n/a & \Res{tsr.test_results.vimeo1080p.plain.3.igpu_s3.temporal_err} & \Res{tsr.test_results.vimeo1080p.plain.3.igpu_s3.tof} & n/a \\
\tier{low} student & \Res{tsr.test_results.vimeo1080p.plain.3.student_s3.psnr_y} & \Res{tsr.test_results.vimeo1080p.plain.3.student_s3.ssim_y} & \Res{tsr.test_results.vimeo1080p.plain.3.student_s3.lpips} & \Res{tsr.test_results.vimeo1080p.plain.3.student_s3.psnr_disocc} & \Res{tsr.test_results.vimeo1080p.plain.3.student_s3.psnr_detail} & n/a & \Res{tsr.test_results.vimeo1080p.plain.3.student_s3.temporal_err} & \Res{tsr.test_results.vimeo1080p.plain.3.student_s3.tof} & n/a \\
Teacher & \Res{tsr.test_results.vimeo1080p.plain.3.teacher_s3.psnr_y} & \Res{tsr.test_results.vimeo1080p.plain.3.teacher_s3.ssim_y} & \Res{tsr.test_results.vimeo1080p.plain.3.teacher_s3.lpips} & \Res{tsr.test_results.vimeo1080p.plain.3.teacher_s3.psnr_disocc} & \Res{tsr.test_results.vimeo1080p.plain.3.teacher_s3.psnr_detail} & n/a & \Res{tsr.test_results.vimeo1080p.plain.3.teacher_s3.temporal_err} & \Res{tsr.test_results.vimeo1080p.plain.3.teacher_s3.tof} & n/a \\
\botrule
\end{tabular*}
\end{table}

\emph{Sintel.} In engine mode the student reaches \Res{tsr.test_results.sintel.engine.3.student_s3.psnr_y} dB, \ResAbsDiff{tsr.test_results.sintel.engine.3.student_s3.psnr_y}{tsr.test_results.sintel.engine.3.analytic_taau.psnr_y} dB above TAA. In plain mode it reaches \Res{tsr.test_results.sintel.plain.3.student_s3.psnr_y} dB, \ResAbsDiff{tsr.test_results.sintel.plain.3.student_s3.psnr_y}{tsr.test_results.sintel.plain.3.bicubic.psnr_y} dB above bicubic. \Preliminary{} The stability measures disagree on this suite. The plain-video temporal-change error of the student is higher than bicubic's (\Res{tsr.test_results.sintel.plain.3.student_s3.temporal_err} against \Res{tsr.test_results.sintel.plain.3.bicubic.temporal_err}), but its tOF is lower (\Res{tsr.test_results.sintel.plain.3.student_s3.tof} against \Res{tsr.test_results.sintel.plain.3.bicubic.tof}). The two measures sample the motion differently, so this suite does not support a claim that the student is more stable in plain mode.

\emph{Vid4.} In engine mode, where the motion comes from RAFT on the clean frames, the student reaches \Res{tsr.test_results.vid4.engine.3.student_s3.psnr_y} dB, \ResAbsDiff{tsr.test_results.vid4.engine.3.student_s3.psnr_y}{tsr.test_results.vid4.engine.3.bicubic.psnr_y} dB above bicubic. In plain mode it reaches \Res{tsr.test_results.vid4.plain.3.student_s3.psnr_y} dB against \Res{tsr.test_results.vid4.plain.3.bicubic.psnr_y} dB, and its LPIPS is \Res{tsr.test_results.vid4.plain.3.student_s3.lpips} against \Res{tsr.test_results.vid4.plain.3.bicubic.lpips}, the largest perceptual gain of the plain suites. \Observation{} \Cref{fig:tsrr-vis-vid4} shows that much of the difference is fine texture, and that the texture the student adds on the cobblestones is a regular pattern, which the ground truth does not share exactly.

\begin{figure}[!htbp]
\centering
\includegraphics[width=\textwidth]{\TsrOut/plots/visual_vid4_plain.png}
\caption{Vid4 as plain video at 3x (bicubic downscale, no motion vectors), last-frame crops, history weight and x-t slices, in the layout of \cref{fig:tsrr-vis-synth}}\label{fig:tsrr-vis-vid4}
\end{figure}

\FloatBarrier

\begin{figure}[!htbp]
\centering
\includegraphics[width=\textwidth]{\TsrOut/plots/visual_synthetic_fast.png}
\caption{Synthetic fast motion at 3x, last-frame crops. Top: bilinear, analytic TAA, \tier{igpu}, \tier{low} student, teacher and ground truth, then the history weight of the student. Bottom: x-t slices of each method and the true disocclusion mask}\label{fig:tsrr-vis-synth}
\end{figure}

\emph{Vimeo-1080p.} In engine mode the student reaches \Res{tsr.test_results.vimeo1080p.engine.3.student_s3.psnr_y} dB, \ResAbsDiff{tsr.test_results.vimeo1080p.engine.3.student_s3.psnr_y}{tsr.test_results.vimeo1080p.engine.3.bicubic.psnr_y} dB above bicubic, with RAFT motion as on Vid4. In plain mode the gain over bicubic is the smallest of all suites, \ResAbsDiff{tsr.test_results.vimeo1080p.plain.3.student_s3.psnr_y}{tsr.test_results.vimeo1080p.plain.3.bicubic.psnr_y} dB, while LPIPS falls from \Res{tsr.test_results.vimeo1080p.plain.3.bicubic.lpips} to \Res{tsr.test_results.vimeo1080p.plain.3.student_s3.lpips}. \Preliminary{} The plain-video temporal-change error of the student is marginally higher than bicubic's (\Res{tsr.test_results.vimeo1080p.plain.3.student_s3.temporal_err} against \Res{tsr.test_results.vimeo1080p.plain.3.bicubic.temporal_err}). Bicubic has the highest plain-mode PSNR-Y of the five suites, \Res{tsr.test_results.vimeo1080p.plain.3.bicubic.psnr_y} dB, which leaves the least room for a learned upscaler. This is a plausible reason for the small gain, and this run does not test it.

\subsubsection{Cross-suite comparisons and the teacher}\label{sec:tsrr-main-summary}

\emph{Plain-video gain over bicubic.} The student beats bicubic in plain mode on all five 3x suites. The mean PSNR-Y gain ranges from \ResAbsDiff{tsr.test_results.vimeo1080p.plain.3.student_s3.psnr_y}{tsr.test_results.vimeo1080p.plain.3.bicubic.psnr_y} dB on Vimeo-1080p to \ResAbsDiff{tsr.test_results.synthetic-game.plain.3.student_s3.psnr_y}{tsr.test_results.synthetic-game.plain.3.bicubic.psnr_y} dB on the synthetic scenes. \Preliminary{} The per-clip counts of the real suites, which are not written as keys, show that the student's clip mean is above bicubic's on every real clip except one, a Vimeo-1080p clip with the highest bicubic score of its suite.

\emph{Engine-mode gain over the analytic TAA baseline.} The student beats TAA on each of the five 3x engine suites. The gap ranges from \ResAbsDiff{tsr.test_results.tartanair.engine.3.student_s3.psnr_y}{tsr.test_results.tartanair.engine.3.analytic_taau.psnr_y} dB on TartanAir to \ResAbsDiff{tsr.test_results.synthetic-game.engine.3.student_s3.psnr_y}{tsr.test_results.synthetic-game.engine.3.analytic_taau.psnr_y} dB on the synthetic scenes. TAA is below bicubic on each of the five engine suites, for example by \ResAbsDiff{tsr.test_results.vimeo1080p.engine.3.bicubic.psnr_y}{tsr.test_results.vimeo1080p.engine.3.analytic_taau.psnr_y} dB on Vimeo-1080p. \Observation{} The baseline is therefore a weak reference, and the engine-mode margins are margins over this reference, not over a production temporal upscaler. A comparison with an engine's own TAA would be needed before these margins are used as a claim about deployed TAA.

\emph{Teacher as the quality ceiling.} The teacher has the same architecture at 64 channels. It costs \Res{tsr.model_tiers.teacher.3.gmac_per_frame_1080p} GMAC per 1080p frame, \ResRatio{tsr.model_tiers.teacher.3.gmac_per_frame_1080p}{tsr.model_tiers.low.3.gmac_per_frame_1080p} times the \tier{low} student, and it is the best 3x model on every suite and mode. In engine mode it exceeds the student by \ResAbsDiff{tsr.test_results.synthetic-game.engine.3.teacher_s3.psnr_y}{tsr.test_results.synthetic-game.engine.3.student_s3.psnr_y} dB on the synthetic scenes and by \ResAbsDiff{tsr.test_results.tartanair.engine.3.teacher_s3.psnr_y}{tsr.test_results.tartanair.engine.3.student_s3.psnr_y} dB on TartanAir. In plain mode the margins are \ResAbsDiff{tsr.test_results.synthetic-game.plain.3.teacher_s3.psnr_y}{tsr.test_results.synthetic-game.plain.3.student_s3.psnr_y} dB and \ResAbsDiff{tsr.test_results.tartanair.plain.3.teacher_s3.psnr_y}{tsr.test_results.tartanair.plain.3.student_s3.psnr_y} dB. \Preliminary{} The teacher is the quality ceiling of this architecture under this budget. Its margin over the student is smaller on the plain suites than on the engine suites, at about ten times the cost, so the teacher is not a deployment candidate.

\emph{The integrated-GPU tier.} The \tier{igpu} tier costs \ResRatio{tsr.model_tiers.igpu.3.gmac_per_frame_1080p}{tsr.model_tiers.low.3.gmac_per_frame_1080p} times the \tier{low} student's MAC count. On the synthetic scenes in plain mode it is \ResAbsDiff{tsr.test_results.synthetic-game.plain.3.igpu_s3.psnr_y}{tsr.test_results.synthetic-game.plain.3.student_s3.psnr_y} dB below the student and \ResAbsDiff{tsr.test_results.synthetic-game.plain.3.igpu_s3.psnr_y}{tsr.test_results.synthetic-game.plain.3.bicubic.psnr_y} dB above bicubic. \Preliminary{} It keeps part of the student's gain at that cost. On Vimeo-1080p in plain mode it scores \Res{tsr.test_results.vimeo1080p.plain.3.igpu_s3.psnr_y} dB, which is \ResAbsDiff{tsr.test_results.vimeo1080p.plain.3.igpu_s3.psnr_y}{tsr.test_results.vimeo1080p.plain.3.bicubic.psnr_y} dB above bicubic and \ResAbsDiff{tsr.test_results.vimeo1080p.plain.3.student_s3.psnr_y}{tsr.test_results.vimeo1080p.plain.3.igpu_s3.psnr_y} dB below the student.

\subsubsection{Per-frame spread}\label{sec:tsrr-perframe}

\Cref{tab:tsrr-perframe} gives the spread of the per-frame PSNR-Y of the student and of bicubic. A clip mean hides the frames where a method fails, and a single bad frame is visible in a real-time renderer. \Preliminary{} The worst frame of the student is above the worst frame of bicubic on every suite and mode in the table. The best frame is a different matter. On Vimeo-1080p the best frame of bicubic is above the best frame of the student in both modes, so the Vimeo gain is not in the sharpest frames. The spread between the best and worst frame is wide for both methods on every suite, because it follows the content of each clip, so the table supports a claim about the worst frames of this run, not about any one clip.

\begin{table}[!htbp]
\caption{Per-frame PSNR-Y (dB) at 3x: minimum, mean and maximum over the per-frame rows of each suite and mode (frames 1 to 15 of each clip), for the \tier{low} student and bicubic. The minimum is the worst frame.}\label{tab:tsrr-perframe}
\footnotesize
\setlength{\tabcolsep}{3pt}
\begin{tabular*}{\textwidth}{@{\extracolsep\fill}llrrrrrr}
\toprule
 & & \multicolumn{3}{c}{\tier{low} student} & \multicolumn{3}{c}{Bicubic} \\
\cmidrule(lr){3-5}\cmidrule(lr){6-8}
Suite & Mode & Min & Mean & Max & Min & Mean & Max \\
\midrule
Synthetic-game & Engine & \Res{tsr.test_per_frame.synthetic-game.engine.3.student_s3.psnr_y.min} & \Res{tsr.test_per_frame.synthetic-game.engine.3.student_s3.psnr_y.mean} & \Res{tsr.test_per_frame.synthetic-game.engine.3.student_s3.psnr_y.max} & \Res{tsr.test_per_frame.synthetic-game.engine.3.bicubic.psnr_y.min} & \Res{tsr.test_per_frame.synthetic-game.engine.3.bicubic.psnr_y.mean} & \Res{tsr.test_per_frame.synthetic-game.engine.3.bicubic.psnr_y.max} \\
Synthetic-game & Plain & \Res{tsr.test_per_frame.synthetic-game.plain.3.student_s3.psnr_y.min} & \Res{tsr.test_per_frame.synthetic-game.plain.3.student_s3.psnr_y.mean} & \Res{tsr.test_per_frame.synthetic-game.plain.3.student_s3.psnr_y.max} & \Res{tsr.test_per_frame.synthetic-game.plain.3.bicubic.psnr_y.min} & \Res{tsr.test_per_frame.synthetic-game.plain.3.bicubic.psnr_y.mean} & \Res{tsr.test_per_frame.synthetic-game.plain.3.bicubic.psnr_y.max} \\
TartanAir & Engine & \Res{tsr.test_per_frame.tartanair.engine.3.student_s3.psnr_y.min} & \Res{tsr.test_per_frame.tartanair.engine.3.student_s3.psnr_y.mean} & \Res{tsr.test_per_frame.tartanair.engine.3.student_s3.psnr_y.max} & \Res{tsr.test_per_frame.tartanair.engine.3.bicubic.psnr_y.min} & \Res{tsr.test_per_frame.tartanair.engine.3.bicubic.psnr_y.mean} & \Res{tsr.test_per_frame.tartanair.engine.3.bicubic.psnr_y.max} \\
TartanAir & Plain & \Res{tsr.test_per_frame.tartanair.plain.3.student_s3.psnr_y.min} & \Res{tsr.test_per_frame.tartanair.plain.3.student_s3.psnr_y.mean} & \Res{tsr.test_per_frame.tartanair.plain.3.student_s3.psnr_y.max} & \Res{tsr.test_per_frame.tartanair.plain.3.bicubic.psnr_y.min} & \Res{tsr.test_per_frame.tartanair.plain.3.bicubic.psnr_y.mean} & \Res{tsr.test_per_frame.tartanair.plain.3.bicubic.psnr_y.max} \\
Sintel & Engine & \Res{tsr.test_per_frame.sintel.engine.3.student_s3.psnr_y.min} & \Res{tsr.test_per_frame.sintel.engine.3.student_s3.psnr_y.mean} & \Res{tsr.test_per_frame.sintel.engine.3.student_s3.psnr_y.max} & \Res{tsr.test_per_frame.sintel.engine.3.bicubic.psnr_y.min} & \Res{tsr.test_per_frame.sintel.engine.3.bicubic.psnr_y.mean} & \Res{tsr.test_per_frame.sintel.engine.3.bicubic.psnr_y.max} \\
Sintel & Plain & \Res{tsr.test_per_frame.sintel.plain.3.student_s3.psnr_y.min} & \Res{tsr.test_per_frame.sintel.plain.3.student_s3.psnr_y.mean} & \Res{tsr.test_per_frame.sintel.plain.3.student_s3.psnr_y.max} & \Res{tsr.test_per_frame.sintel.plain.3.bicubic.psnr_y.min} & \Res{tsr.test_per_frame.sintel.plain.3.bicubic.psnr_y.mean} & \Res{tsr.test_per_frame.sintel.plain.3.bicubic.psnr_y.max} \\
Vid4 & Engine & \Res{tsr.test_per_frame.vid4.engine.3.student_s3.psnr_y.min} & \Res{tsr.test_per_frame.vid4.engine.3.student_s3.psnr_y.mean} & \Res{tsr.test_per_frame.vid4.engine.3.student_s3.psnr_y.max} & \Res{tsr.test_per_frame.vid4.engine.3.bicubic.psnr_y.min} & \Res{tsr.test_per_frame.vid4.engine.3.bicubic.psnr_y.mean} & \Res{tsr.test_per_frame.vid4.engine.3.bicubic.psnr_y.max} \\
Vid4 & Plain & \Res{tsr.test_per_frame.vid4.plain.3.student_s3.psnr_y.min} & \Res{tsr.test_per_frame.vid4.plain.3.student_s3.psnr_y.mean} & \Res{tsr.test_per_frame.vid4.plain.3.student_s3.psnr_y.max} & \Res{tsr.test_per_frame.vid4.plain.3.bicubic.psnr_y.min} & \Res{tsr.test_per_frame.vid4.plain.3.bicubic.psnr_y.mean} & \Res{tsr.test_per_frame.vid4.plain.3.bicubic.psnr_y.max} \\
Vimeo-1080p & Engine & \Res{tsr.test_per_frame.vimeo1080p.engine.3.student_s3.psnr_y.min} & \Res{tsr.test_per_frame.vimeo1080p.engine.3.student_s3.psnr_y.mean} & \Res{tsr.test_per_frame.vimeo1080p.engine.3.student_s3.psnr_y.max} & \Res{tsr.test_per_frame.vimeo1080p.engine.3.bicubic.psnr_y.min} & \Res{tsr.test_per_frame.vimeo1080p.engine.3.bicubic.psnr_y.mean} & \Res{tsr.test_per_frame.vimeo1080p.engine.3.bicubic.psnr_y.max} \\
Vimeo-1080p & Plain & \Res{tsr.test_per_frame.vimeo1080p.plain.3.student_s3.psnr_y.min} & \Res{tsr.test_per_frame.vimeo1080p.plain.3.student_s3.psnr_y.mean} & \Res{tsr.test_per_frame.vimeo1080p.plain.3.student_s3.psnr_y.max} & \Res{tsr.test_per_frame.vimeo1080p.plain.3.bicubic.psnr_y.min} & \Res{tsr.test_per_frame.vimeo1080p.plain.3.bicubic.psnr_y.mean} & \Res{tsr.test_per_frame.vimeo1080p.plain.3.bicubic.psnr_y.max} \\
\botrule
\end{tabular*}
\end{table}

\subsection{Results at 2x and on native renders}\label{sec:tsrr-2x}

\footnotesize
\setlength{\tabcolsep}{2.5pt}
\begin{longtable}{@{}lrrrrrrrrr@{}}
\caption{Results at 2x: the \tier{low} 2x student against bicubic, bilinear and analytic TAA (TAA in engine mode only). Means over per-frame rows. Disocc., Detail and Part.: PSNR-Y inside disoccluded, fine-detail and particle pixels. Temp.: temporal-change error. Warp: warping error. tOF: temporal optical-flow error in pixels. n/a: not computed for this suite.}\label{tab:tsrr-2x}\\
\toprule
Method & PSNR-Y & SSIM-Y & LPIPS & Disocc. & Detail & Part. & Temp. & tOF & Warp \\
\midrule
\endfirsthead
\multicolumn{10}{l}{\emph{Table \thetable, continued}} \\
\toprule
Method & PSNR-Y & SSIM-Y & LPIPS & Disocc. & Detail & Part. & Temp. & tOF & Warp \\
\midrule
\endhead
\midrule
\multicolumn{10}{r}{\emph{Continued on next page}} \\
\endfoot
\botrule
\endlastfoot
\multicolumn{10}{@{}l}{\emph{Synthetic-game, engine mode}} \\
Bicubic & \Res{tsr.test_results.synthetic-game.engine.2.bicubic.psnr_y} & \Res{tsr.test_results.synthetic-game.engine.2.bicubic.ssim_y} & \Res{tsr.test_results.synthetic-game.engine.2.bicubic.lpips} & \Res{tsr.test_results.synthetic-game.engine.2.bicubic.psnr_disocc} & \Res{tsr.test_results.synthetic-game.engine.2.bicubic.psnr_detail} & \Res{tsr.test_results.synthetic-game.engine.2.bicubic.psnr_particles} & \Res{tsr.test_results.synthetic-game.engine.2.bicubic.temporal_err} & \Res{tsr.test_results.synthetic-game.engine.2.bicubic.tof} & \Res{tsr.test_results.synthetic-game.engine.2.bicubic.warp_err} \\
Bilinear & \Res{tsr.test_results.synthetic-game.engine.2.bilinear.psnr_y} & \Res{tsr.test_results.synthetic-game.engine.2.bilinear.ssim_y} & \Res{tsr.test_results.synthetic-game.engine.2.bilinear.lpips} & \Res{tsr.test_results.synthetic-game.engine.2.bilinear.psnr_disocc} & \Res{tsr.test_results.synthetic-game.engine.2.bilinear.psnr_detail} & \Res{tsr.test_results.synthetic-game.engine.2.bilinear.psnr_particles} & \Res{tsr.test_results.synthetic-game.engine.2.bilinear.temporal_err} & \Res{tsr.test_results.synthetic-game.engine.2.bilinear.tof} & \Res{tsr.test_results.synthetic-game.engine.2.bilinear.warp_err} \\
Analytic TAA & \Res{tsr.test_results.synthetic-game.engine.2.analytic_taau.psnr_y} & \Res{tsr.test_results.synthetic-game.engine.2.analytic_taau.ssim_y} & \Res{tsr.test_results.synthetic-game.engine.2.analytic_taau.lpips} & \Res{tsr.test_results.synthetic-game.engine.2.analytic_taau.psnr_disocc} & \Res{tsr.test_results.synthetic-game.engine.2.analytic_taau.psnr_detail} & \Res{tsr.test_results.synthetic-game.engine.2.analytic_taau.psnr_particles} & \Res{tsr.test_results.synthetic-game.engine.2.analytic_taau.temporal_err} & \Res{tsr.test_results.synthetic-game.engine.2.analytic_taau.tof} & \Res{tsr.test_results.synthetic-game.engine.2.analytic_taau.warp_err} \\
\tier{low} 2x student & \Res{tsr.test_results.synthetic-game.engine.2.student_s2.psnr_y} & \Res{tsr.test_results.synthetic-game.engine.2.student_s2.ssim_y} & \Res{tsr.test_results.synthetic-game.engine.2.student_s2.lpips} & \Res{tsr.test_results.synthetic-game.engine.2.student_s2.psnr_disocc} & \Res{tsr.test_results.synthetic-game.engine.2.student_s2.psnr_detail} & \Res{tsr.test_results.synthetic-game.engine.2.student_s2.psnr_particles} & \Res{tsr.test_results.synthetic-game.engine.2.student_s2.temporal_err} & \Res{tsr.test_results.synthetic-game.engine.2.student_s2.tof} & \Res{tsr.test_results.synthetic-game.engine.2.student_s2.warp_err} \\
\multicolumn{10}{@{}l}{\emph{Synthetic-game, plain mode}} \\
Bicubic & \Res{tsr.test_results.synthetic-game.plain.2.bicubic.psnr_y} & \Res{tsr.test_results.synthetic-game.plain.2.bicubic.ssim_y} & \Res{tsr.test_results.synthetic-game.plain.2.bicubic.lpips} & \Res{tsr.test_results.synthetic-game.plain.2.bicubic.psnr_disocc} & \Res{tsr.test_results.synthetic-game.plain.2.bicubic.psnr_detail} & n/a & \Res{tsr.test_results.synthetic-game.plain.2.bicubic.temporal_err} & \Res{tsr.test_results.synthetic-game.plain.2.bicubic.tof} & \Res{tsr.test_results.synthetic-game.plain.2.bicubic.warp_err} \\
Bilinear & \Res{tsr.test_results.synthetic-game.plain.2.bilinear.psnr_y} & \Res{tsr.test_results.synthetic-game.plain.2.bilinear.ssim_y} & \Res{tsr.test_results.synthetic-game.plain.2.bilinear.lpips} & \Res{tsr.test_results.synthetic-game.plain.2.bilinear.psnr_disocc} & \Res{tsr.test_results.synthetic-game.plain.2.bilinear.psnr_detail} & n/a & \Res{tsr.test_results.synthetic-game.plain.2.bilinear.temporal_err} & \Res{tsr.test_results.synthetic-game.plain.2.bilinear.tof} & \Res{tsr.test_results.synthetic-game.plain.2.bilinear.warp_err} \\
\tier{low} 2x student & \Res{tsr.test_results.synthetic-game.plain.2.student_s2.psnr_y} & \Res{tsr.test_results.synthetic-game.plain.2.student_s2.ssim_y} & \Res{tsr.test_results.synthetic-game.plain.2.student_s2.lpips} & \Res{tsr.test_results.synthetic-game.plain.2.student_s2.psnr_disocc} & \Res{tsr.test_results.synthetic-game.plain.2.student_s2.psnr_detail} & n/a & \Res{tsr.test_results.synthetic-game.plain.2.student_s2.temporal_err} & \Res{tsr.test_results.synthetic-game.plain.2.student_s2.tof} & \Res{tsr.test_results.synthetic-game.plain.2.student_s2.warp_err} \\
\multicolumn{10}{@{}l}{\emph{TartanAir, engine mode}} \\
Bicubic & \Res{tsr.test_results.tartanair.engine.2.bicubic.psnr_y} & \Res{tsr.test_results.tartanair.engine.2.bicubic.ssim_y} & \Res{tsr.test_results.tartanair.engine.2.bicubic.lpips} & \Res{tsr.test_results.tartanair.engine.2.bicubic.psnr_disocc} & \Res{tsr.test_results.tartanair.engine.2.bicubic.psnr_detail} & n/a & \Res{tsr.test_results.tartanair.engine.2.bicubic.temporal_err} & \Res{tsr.test_results.tartanair.engine.2.bicubic.tof} & \Res{tsr.test_results.tartanair.engine.2.bicubic.warp_err} \\
Bilinear & \Res{tsr.test_results.tartanair.engine.2.bilinear.psnr_y} & \Res{tsr.test_results.tartanair.engine.2.bilinear.ssim_y} & \Res{tsr.test_results.tartanair.engine.2.bilinear.lpips} & \Res{tsr.test_results.tartanair.engine.2.bilinear.psnr_disocc} & \Res{tsr.test_results.tartanair.engine.2.bilinear.psnr_detail} & n/a & \Res{tsr.test_results.tartanair.engine.2.bilinear.temporal_err} & \Res{tsr.test_results.tartanair.engine.2.bilinear.tof} & \Res{tsr.test_results.tartanair.engine.2.bilinear.warp_err} \\
Analytic TAA & \Res{tsr.test_results.tartanair.engine.2.analytic_taau.psnr_y} & \Res{tsr.test_results.tartanair.engine.2.analytic_taau.ssim_y} & \Res{tsr.test_results.tartanair.engine.2.analytic_taau.lpips} & \Res{tsr.test_results.tartanair.engine.2.analytic_taau.psnr_disocc} & \Res{tsr.test_results.tartanair.engine.2.analytic_taau.psnr_detail} & n/a & \Res{tsr.test_results.tartanair.engine.2.analytic_taau.temporal_err} & \Res{tsr.test_results.tartanair.engine.2.analytic_taau.tof} & \Res{tsr.test_results.tartanair.engine.2.analytic_taau.warp_err} \\
\tier{low} 2x student & \Res{tsr.test_results.tartanair.engine.2.student_s2.psnr_y} & \Res{tsr.test_results.tartanair.engine.2.student_s2.ssim_y} & \Res{tsr.test_results.tartanair.engine.2.student_s2.lpips} & \Res{tsr.test_results.tartanair.engine.2.student_s2.psnr_disocc} & \Res{tsr.test_results.tartanair.engine.2.student_s2.psnr_detail} & n/a & \Res{tsr.test_results.tartanair.engine.2.student_s2.temporal_err} & \Res{tsr.test_results.tartanair.engine.2.student_s2.tof} & \Res{tsr.test_results.tartanair.engine.2.student_s2.warp_err} \\
\multicolumn{10}{@{}l}{\emph{TartanAir, plain mode}} \\
Bicubic & \Res{tsr.test_results.tartanair.plain.2.bicubic.psnr_y} & \Res{tsr.test_results.tartanair.plain.2.bicubic.ssim_y} & \Res{tsr.test_results.tartanair.plain.2.bicubic.lpips} & \Res{tsr.test_results.tartanair.plain.2.bicubic.psnr_disocc} & \Res{tsr.test_results.tartanair.plain.2.bicubic.psnr_detail} & n/a & \Res{tsr.test_results.tartanair.plain.2.bicubic.temporal_err} & \Res{tsr.test_results.tartanair.plain.2.bicubic.tof} & \Res{tsr.test_results.tartanair.plain.2.bicubic.warp_err} \\
Bilinear & \Res{tsr.test_results.tartanair.plain.2.bilinear.psnr_y} & \Res{tsr.test_results.tartanair.plain.2.bilinear.ssim_y} & \Res{tsr.test_results.tartanair.plain.2.bilinear.lpips} & \Res{tsr.test_results.tartanair.plain.2.bilinear.psnr_disocc} & \Res{tsr.test_results.tartanair.plain.2.bilinear.psnr_detail} & n/a & \Res{tsr.test_results.tartanair.plain.2.bilinear.temporal_err} & \Res{tsr.test_results.tartanair.plain.2.bilinear.tof} & \Res{tsr.test_results.tartanair.plain.2.bilinear.warp_err} \\
\tier{low} 2x student & \Res{tsr.test_results.tartanair.plain.2.student_s2.psnr_y} & \Res{tsr.test_results.tartanair.plain.2.student_s2.ssim_y} & \Res{tsr.test_results.tartanair.plain.2.student_s2.lpips} & \Res{tsr.test_results.tartanair.plain.2.student_s2.psnr_disocc} & \Res{tsr.test_results.tartanair.plain.2.student_s2.psnr_detail} & n/a & \Res{tsr.test_results.tartanair.plain.2.student_s2.temporal_err} & \Res{tsr.test_results.tartanair.plain.2.student_s2.tof} & \Res{tsr.test_results.tartanair.plain.2.student_s2.warp_err} \\
\multicolumn{10}{@{}l}{\emph{Sintel, engine mode}} \\
Bicubic & \Res{tsr.test_results.sintel.engine.2.bicubic.psnr_y} & \Res{tsr.test_results.sintel.engine.2.bicubic.ssim_y} & \Res{tsr.test_results.sintel.engine.2.bicubic.lpips} & \Res{tsr.test_results.sintel.engine.2.bicubic.psnr_disocc} & \Res{tsr.test_results.sintel.engine.2.bicubic.psnr_detail} & n/a & \Res{tsr.test_results.sintel.engine.2.bicubic.temporal_err} & \Res{tsr.test_results.sintel.engine.2.bicubic.tof} & \Res{tsr.test_results.sintel.engine.2.bicubic.warp_err} \\
Bilinear & \Res{tsr.test_results.sintel.engine.2.bilinear.psnr_y} & \Res{tsr.test_results.sintel.engine.2.bilinear.ssim_y} & \Res{tsr.test_results.sintel.engine.2.bilinear.lpips} & \Res{tsr.test_results.sintel.engine.2.bilinear.psnr_disocc} & \Res{tsr.test_results.sintel.engine.2.bilinear.psnr_detail} & n/a & \Res{tsr.test_results.sintel.engine.2.bilinear.temporal_err} & \Res{tsr.test_results.sintel.engine.2.bilinear.tof} & \Res{tsr.test_results.sintel.engine.2.bilinear.warp_err} \\
Analytic TAA & \Res{tsr.test_results.sintel.engine.2.analytic_taau.psnr_y} & \Res{tsr.test_results.sintel.engine.2.analytic_taau.ssim_y} & \Res{tsr.test_results.sintel.engine.2.analytic_taau.lpips} & \Res{tsr.test_results.sintel.engine.2.analytic_taau.psnr_disocc} & \Res{tsr.test_results.sintel.engine.2.analytic_taau.psnr_detail} & n/a & \Res{tsr.test_results.sintel.engine.2.analytic_taau.temporal_err} & \Res{tsr.test_results.sintel.engine.2.analytic_taau.tof} & \Res{tsr.test_results.sintel.engine.2.analytic_taau.warp_err} \\
\tier{low} 2x student & \Res{tsr.test_results.sintel.engine.2.student_s2.psnr_y} & \Res{tsr.test_results.sintel.engine.2.student_s2.ssim_y} & \Res{tsr.test_results.sintel.engine.2.student_s2.lpips} & \Res{tsr.test_results.sintel.engine.2.student_s2.psnr_disocc} & \Res{tsr.test_results.sintel.engine.2.student_s2.psnr_detail} & n/a & \Res{tsr.test_results.sintel.engine.2.student_s2.temporal_err} & \Res{tsr.test_results.sintel.engine.2.student_s2.tof} & \Res{tsr.test_results.sintel.engine.2.student_s2.warp_err} \\
\multicolumn{10}{@{}l}{\emph{Sintel, plain mode}} \\
Bicubic & \Res{tsr.test_results.sintel.plain.2.bicubic.psnr_y} & \Res{tsr.test_results.sintel.plain.2.bicubic.ssim_y} & \Res{tsr.test_results.sintel.plain.2.bicubic.lpips} & \Res{tsr.test_results.sintel.plain.2.bicubic.psnr_disocc} & \Res{tsr.test_results.sintel.plain.2.bicubic.psnr_detail} & n/a & \Res{tsr.test_results.sintel.plain.2.bicubic.temporal_err} & \Res{tsr.test_results.sintel.plain.2.bicubic.tof} & \Res{tsr.test_results.sintel.plain.2.bicubic.warp_err} \\
Bilinear & \Res{tsr.test_results.sintel.plain.2.bilinear.psnr_y} & \Res{tsr.test_results.sintel.plain.2.bilinear.ssim_y} & \Res{tsr.test_results.sintel.plain.2.bilinear.lpips} & \Res{tsr.test_results.sintel.plain.2.bilinear.psnr_disocc} & \Res{tsr.test_results.sintel.plain.2.bilinear.psnr_detail} & n/a & \Res{tsr.test_results.sintel.plain.2.bilinear.temporal_err} & \Res{tsr.test_results.sintel.plain.2.bilinear.tof} & \Res{tsr.test_results.sintel.plain.2.bilinear.warp_err} \\
\tier{low} 2x student & \Res{tsr.test_results.sintel.plain.2.student_s2.psnr_y} & \Res{tsr.test_results.sintel.plain.2.student_s2.ssim_y} & \Res{tsr.test_results.sintel.plain.2.student_s2.lpips} & \Res{tsr.test_results.sintel.plain.2.student_s2.psnr_disocc} & \Res{tsr.test_results.sintel.plain.2.student_s2.psnr_detail} & n/a & \Res{tsr.test_results.sintel.plain.2.student_s2.temporal_err} & \Res{tsr.test_results.sintel.plain.2.student_s2.tof} & \Res{tsr.test_results.sintel.plain.2.student_s2.warp_err} \\
\multicolumn{10}{@{}l}{\emph{Vid4, engine mode}} \\
Bicubic & \Res{tsr.test_results.vid4.engine.2.bicubic.psnr_y} & \Res{tsr.test_results.vid4.engine.2.bicubic.ssim_y} & \Res{tsr.test_results.vid4.engine.2.bicubic.lpips} & \Res{tsr.test_results.vid4.engine.2.bicubic.psnr_disocc} & \Res{tsr.test_results.vid4.engine.2.bicubic.psnr_detail} & n/a & \Res{tsr.test_results.vid4.engine.2.bicubic.temporal_err} & \Res{tsr.test_results.vid4.engine.2.bicubic.tof} & n/a \\
Bilinear & \Res{tsr.test_results.vid4.engine.2.bilinear.psnr_y} & \Res{tsr.test_results.vid4.engine.2.bilinear.ssim_y} & \Res{tsr.test_results.vid4.engine.2.bilinear.lpips} & \Res{tsr.test_results.vid4.engine.2.bilinear.psnr_disocc} & \Res{tsr.test_results.vid4.engine.2.bilinear.psnr_detail} & n/a & \Res{tsr.test_results.vid4.engine.2.bilinear.temporal_err} & \Res{tsr.test_results.vid4.engine.2.bilinear.tof} & n/a \\
Analytic TAA & \Res{tsr.test_results.vid4.engine.2.analytic_taau.psnr_y} & \Res{tsr.test_results.vid4.engine.2.analytic_taau.ssim_y} & \Res{tsr.test_results.vid4.engine.2.analytic_taau.lpips} & \Res{tsr.test_results.vid4.engine.2.analytic_taau.psnr_disocc} & \Res{tsr.test_results.vid4.engine.2.analytic_taau.psnr_detail} & n/a & \Res{tsr.test_results.vid4.engine.2.analytic_taau.temporal_err} & \Res{tsr.test_results.vid4.engine.2.analytic_taau.tof} & n/a \\
\tier{low} 2x student & \Res{tsr.test_results.vid4.engine.2.student_s2.psnr_y} & \Res{tsr.test_results.vid4.engine.2.student_s2.ssim_y} & \Res{tsr.test_results.vid4.engine.2.student_s2.lpips} & \Res{tsr.test_results.vid4.engine.2.student_s2.psnr_disocc} & \Res{tsr.test_results.vid4.engine.2.student_s2.psnr_detail} & n/a & \Res{tsr.test_results.vid4.engine.2.student_s2.temporal_err} & \Res{tsr.test_results.vid4.engine.2.student_s2.tof} & n/a \\
\multicolumn{10}{@{}l}{\emph{Vid4, plain mode}} \\
Bicubic & \Res{tsr.test_results.vid4.plain.2.bicubic.psnr_y} & \Res{tsr.test_results.vid4.plain.2.bicubic.ssim_y} & \Res{tsr.test_results.vid4.plain.2.bicubic.lpips} & \Res{tsr.test_results.vid4.plain.2.bicubic.psnr_disocc} & \Res{tsr.test_results.vid4.plain.2.bicubic.psnr_detail} & n/a & \Res{tsr.test_results.vid4.plain.2.bicubic.temporal_err} & \Res{tsr.test_results.vid4.plain.2.bicubic.tof} & n/a \\
Bilinear & \Res{tsr.test_results.vid4.plain.2.bilinear.psnr_y} & \Res{tsr.test_results.vid4.plain.2.bilinear.ssim_y} & \Res{tsr.test_results.vid4.plain.2.bilinear.lpips} & \Res{tsr.test_results.vid4.plain.2.bilinear.psnr_disocc} & \Res{tsr.test_results.vid4.plain.2.bilinear.psnr_detail} & n/a & \Res{tsr.test_results.vid4.plain.2.bilinear.temporal_err} & \Res{tsr.test_results.vid4.plain.2.bilinear.tof} & n/a \\
\tier{low} 2x student & \Res{tsr.test_results.vid4.plain.2.student_s2.psnr_y} & \Res{tsr.test_results.vid4.plain.2.student_s2.ssim_y} & \Res{tsr.test_results.vid4.plain.2.student_s2.lpips} & \Res{tsr.test_results.vid4.plain.2.student_s2.psnr_disocc} & \Res{tsr.test_results.vid4.plain.2.student_s2.psnr_detail} & n/a & \Res{tsr.test_results.vid4.plain.2.student_s2.temporal_err} & \Res{tsr.test_results.vid4.plain.2.student_s2.tof} & n/a \\
\multicolumn{10}{@{}l}{\emph{Vimeo-1080p, engine mode}} \\
Bicubic & \Res{tsr.test_results.vimeo1080p.engine.2.bicubic.psnr_y} & \Res{tsr.test_results.vimeo1080p.engine.2.bicubic.ssim_y} & \Res{tsr.test_results.vimeo1080p.engine.2.bicubic.lpips} & \Res{tsr.test_results.vimeo1080p.engine.2.bicubic.psnr_disocc} & \Res{tsr.test_results.vimeo1080p.engine.2.bicubic.psnr_detail} & n/a & \Res{tsr.test_results.vimeo1080p.engine.2.bicubic.temporal_err} & \Res{tsr.test_results.vimeo1080p.engine.2.bicubic.tof} & n/a \\
Bilinear & \Res{tsr.test_results.vimeo1080p.engine.2.bilinear.psnr_y} & \Res{tsr.test_results.vimeo1080p.engine.2.bilinear.ssim_y} & \Res{tsr.test_results.vimeo1080p.engine.2.bilinear.lpips} & \Res{tsr.test_results.vimeo1080p.engine.2.bilinear.psnr_disocc} & \Res{tsr.test_results.vimeo1080p.engine.2.bilinear.psnr_detail} & n/a & \Res{tsr.test_results.vimeo1080p.engine.2.bilinear.temporal_err} & \Res{tsr.test_results.vimeo1080p.engine.2.bilinear.tof} & n/a \\
Analytic TAA & \Res{tsr.test_results.vimeo1080p.engine.2.analytic_taau.psnr_y} & \Res{tsr.test_results.vimeo1080p.engine.2.analytic_taau.ssim_y} & \Res{tsr.test_results.vimeo1080p.engine.2.analytic_taau.lpips} & \Res{tsr.test_results.vimeo1080p.engine.2.analytic_taau.psnr_disocc} & \Res{tsr.test_results.vimeo1080p.engine.2.analytic_taau.psnr_detail} & n/a & \Res{tsr.test_results.vimeo1080p.engine.2.analytic_taau.temporal_err} & \Res{tsr.test_results.vimeo1080p.engine.2.analytic_taau.tof} & n/a \\
\tier{low} 2x student & \Res{tsr.test_results.vimeo1080p.engine.2.student_s2.psnr_y} & \Res{tsr.test_results.vimeo1080p.engine.2.student_s2.ssim_y} & \Res{tsr.test_results.vimeo1080p.engine.2.student_s2.lpips} & \Res{tsr.test_results.vimeo1080p.engine.2.student_s2.psnr_disocc} & \Res{tsr.test_results.vimeo1080p.engine.2.student_s2.psnr_detail} & n/a & \Res{tsr.test_results.vimeo1080p.engine.2.student_s2.temporal_err} & \Res{tsr.test_results.vimeo1080p.engine.2.student_s2.tof} & n/a \\
\multicolumn{10}{@{}l}{\emph{Vimeo-1080p, plain mode}} \\
Bicubic & \Res{tsr.test_results.vimeo1080p.plain.2.bicubic.psnr_y} & \Res{tsr.test_results.vimeo1080p.plain.2.bicubic.ssim_y} & \Res{tsr.test_results.vimeo1080p.plain.2.bicubic.lpips} & \Res{tsr.test_results.vimeo1080p.plain.2.bicubic.psnr_disocc} & \Res{tsr.test_results.vimeo1080p.plain.2.bicubic.psnr_detail} & n/a & \Res{tsr.test_results.vimeo1080p.plain.2.bicubic.temporal_err} & \Res{tsr.test_results.vimeo1080p.plain.2.bicubic.tof} & n/a \\
Bilinear & \Res{tsr.test_results.vimeo1080p.plain.2.bilinear.psnr_y} & \Res{tsr.test_results.vimeo1080p.plain.2.bilinear.ssim_y} & \Res{tsr.test_results.vimeo1080p.plain.2.bilinear.lpips} & \Res{tsr.test_results.vimeo1080p.plain.2.bilinear.psnr_disocc} & \Res{tsr.test_results.vimeo1080p.plain.2.bilinear.psnr_detail} & n/a & \Res{tsr.test_results.vimeo1080p.plain.2.bilinear.temporal_err} & \Res{tsr.test_results.vimeo1080p.plain.2.bilinear.tof} & n/a \\
\tier{low} 2x student & \Res{tsr.test_results.vimeo1080p.plain.2.student_s2.psnr_y} & \Res{tsr.test_results.vimeo1080p.plain.2.student_s2.ssim_y} & \Res{tsr.test_results.vimeo1080p.plain.2.student_s2.lpips} & \Res{tsr.test_results.vimeo1080p.plain.2.student_s2.psnr_disocc} & \Res{tsr.test_results.vimeo1080p.plain.2.student_s2.psnr_detail} & n/a & \Res{tsr.test_results.vimeo1080p.plain.2.student_s2.temporal_err} & \Res{tsr.test_results.vimeo1080p.plain.2.student_s2.tof} & n/a \\
\multicolumn{10}{@{}l}{\emph{GameIR native, engine mode}} \\
Bicubic & \Res{tsr.test_results.gameir-native.engine.2.bicubic.psnr_y} & \Res{tsr.test_results.gameir-native.engine.2.bicubic.ssim_y} & \Res{tsr.test_results.gameir-native.engine.2.bicubic.lpips} & \Res{tsr.test_results.gameir-native.engine.2.bicubic.psnr_disocc} & \Res{tsr.test_results.gameir-native.engine.2.bicubic.psnr_detail} & n/a & \Res{tsr.test_results.gameir-native.engine.2.bicubic.temporal_err} & \Res{tsr.test_results.gameir-native.engine.2.bicubic.tof} & n/a \\
Bilinear & \Res{tsr.test_results.gameir-native.engine.2.bilinear.psnr_y} & \Res{tsr.test_results.gameir-native.engine.2.bilinear.ssim_y} & \Res{tsr.test_results.gameir-native.engine.2.bilinear.lpips} & \Res{tsr.test_results.gameir-native.engine.2.bilinear.psnr_disocc} & \Res{tsr.test_results.gameir-native.engine.2.bilinear.psnr_detail} & n/a & \Res{tsr.test_results.gameir-native.engine.2.bilinear.temporal_err} & \Res{tsr.test_results.gameir-native.engine.2.bilinear.tof} & n/a \\
Analytic TAA & \Res{tsr.test_results.gameir-native.engine.2.analytic_taau.psnr_y} & \Res{tsr.test_results.gameir-native.engine.2.analytic_taau.ssim_y} & \Res{tsr.test_results.gameir-native.engine.2.analytic_taau.lpips} & \Res{tsr.test_results.gameir-native.engine.2.analytic_taau.psnr_disocc} & \Res{tsr.test_results.gameir-native.engine.2.analytic_taau.psnr_detail} & n/a & \Res{tsr.test_results.gameir-native.engine.2.analytic_taau.temporal_err} & \Res{tsr.test_results.gameir-native.engine.2.analytic_taau.tof} & n/a \\
\tier{low} 2x student & \Res{tsr.test_results.gameir-native.engine.2.student_s2.psnr_y} & \Res{tsr.test_results.gameir-native.engine.2.student_s2.ssim_y} & \Res{tsr.test_results.gameir-native.engine.2.student_s2.lpips} & \Res{tsr.test_results.gameir-native.engine.2.student_s2.psnr_disocc} & \Res{tsr.test_results.gameir-native.engine.2.student_s2.psnr_detail} & n/a & \Res{tsr.test_results.gameir-native.engine.2.student_s2.temporal_err} & \Res{tsr.test_results.gameir-native.engine.2.student_s2.tof} & n/a \\
\end{longtable}

\Cref{tab:tsrr-2x} gives the 2x results. The \tier{low} 2x student is compared with bicubic, bilinear and TAA, and the GameIR native suite is run in engine mode only. In plain mode the 2x student beats bicubic on every suite, by \ResAbsDiff{tsr.test_results.vimeo1080p.plain.2.student_s2.psnr_y}{tsr.test_results.vimeo1080p.plain.2.bicubic.psnr_y} dB on Vimeo-1080p to \ResAbsDiff{tsr.test_results.synthetic-game.plain.2.student_s2.psnr_y}{tsr.test_results.synthetic-game.plain.2.bicubic.psnr_y} dB on the synthetic scenes. \Preliminary{} The plain-mode gain at 2x is larger than at 3x on every suite. A 2x input keeps more of the target's detail, which is a plausible reason, but the notebook does not test it.

\emph{GameIR native.} The student reaches \Res{tsr.test_results.gameir-native.engine.2.student_s2.psnr_y} dB on the native renders, against \Res{tsr.test_results.gameir-native.engine.2.bicubic.psnr_y} dB for bicubic and \Res{tsr.test_results.gameir-native.engine.2.analytic_taau.psnr_y} dB for TAA. The gain over bicubic is \ResAbsDiff{tsr.test_results.gameir-native.engine.2.student_s2.psnr_y}{tsr.test_results.gameir-native.engine.2.bicubic.psnr_y} dB, which is small. Its LPIPS is lower (\Res{tsr.test_results.gameir-native.engine.2.student_s2.lpips} against \Res{tsr.test_results.gameir-native.engine.2.bicubic.lpips}), but its tOF is higher (\Res{tsr.test_results.gameir-native.engine.2.student_s2.tof} against \Res{tsr.test_results.gameir-native.engine.2.bicubic.tof}). \Preliminary{} This is the only real-renderer test, and it is the test closest to the intended input. It rests on four clips, one LPIPS frame per clip and RAFT motion from the clean frames. It supports the claim that the 2x student does not lose to bicubic on native renders, but it does not establish the size of the gain or any claim about stability.

\subsection{Quality against stability}\label{sec:tsrr-stability}

\Cref{fig:tsrr-stability} plots each method's PSNR-Y against its temporal-change error for 3x engine mode (left) and plain mode (middle), and gives the plain-video PSNR-Y of each suite (right). A better method lies towards the top left, because it is sharper and more stable. The points of one suite lie close together, so the scatter panels compare methods mostly within a suite. The labels overlap in the scatter panels, so the tables of \cref{sec:tsrr-main} are the reference for any single point. \Observation{} In engine mode the trained tiers and the teacher form a cluster at the top left. Bicubic and bilinear lie towards the right, which means less stable, and the analytic TAA lies low on the PSNR axis, which means stable but blurred. In plain mode the methods differ little in temporal-change error within a suite, and the trained models separate from the baselines mainly in PSNR.

\subsection{Motion regions and hard regions}\label{sec:tsrr-motion}

\Cref{fig:tsrr-motion} splits the synthetic engine results by motion class (left), gives the PSNR-Y in the three hard regions (middle), and gives LPIPS and tOF (right). The hard-region values for the main methods are in \cref{tab:tsrr-main-synthetic-game}. In every hard region the student is ahead of TAA and bicubic, and the teacher is ahead of the student. \Preliminary{} The class means come from the per-frame file, which the generator does not yet aggregate by class, so the left panel is read qualitatively. \Observation{} The class means are not a function of speed alone, because each class uses its own test textures. The fast class has the highest PSNR-Y for every method in the left panel, so the class differences mix speed and content. Within each class the student and the teacher are above the baselines. The gap to the baselines is largest in the static class in engine mode, where the history accumulates without the disturbance of motion.

\begin{figure}[!htbp]
\centering
\includegraphics[width=\textwidth]{\TsrOut/plots/motion_regions.png}
\caption{Synthetic-game at 3x in engine mode. Left: PSNR-Y by motion class. Middle: PSNR-Y in disoccluded, fine-detail and particle regions. Right: LPIPS and tOF. The bars are means over the test stills, and lower is better for LPIPS and tOF}\label{fig:tsrr-motion}
\end{figure}

\emph{Disocclusions.} A disoccluded pixel has no valid history, so a temporal network must reject the history there. The student's disocclusion PSNR is \ResAbsDiff{tsr.test_results.synthetic-game.engine.3.student_s3.psnr_disocc}{tsr.test_results.synthetic-game.engine.3.bicubic.psnr_disocc} dB above bicubic, a smaller gap than its \ResAbsDiff{tsr.test_results.synthetic-game.engine.3.student_s3.psnr_detail}{tsr.test_results.synthetic-game.engine.3.bicubic.psnr_detail} dB gap in the fine-detail region. \Observation{} The history weight in \cref{fig:tsrr-vis-synth} shows the mechanism, because it drops to black along the leading edge of the moving object, where the true disocclusion lies.

\subsection{Static-scene convergence}\label{sec:tsrr-static}

With a static camera and sub-pixel jitter, each frame adds new sub-pixel samples, so a temporal network can in principle reconstruct the native image as the sequence continues. \Cref{fig:tsrr-static} and \cref{tab:tsrr-static} show PSNR-Y against frame index for the static-camera methods. Single-frame and fixed-rule methods stay flat. \Preliminary{} The student rises over the 16 frames, and the teacher rises faster and to a higher level. The 30-minute reference rises to about frame 6 and then declines slightly. Without the 16-frame unrolls the curve peaks earlier and falls well below its peak by frame 15, and its last-frame value is the lowest of the recurrent models. The single-frame ablation stays near the level of TAA, which means that the gain over frame 0 comes from the history and not from the network's single-frame upscaling. This test uses engine input only, so plain-video convergence was not measured.

\begin{figure}[!htbp]
\centering
\includegraphics[width=0.85\textwidth,height=0.25\textheight,keepaspectratio]{\TsrOut/plots/static_convergence.png}
\caption{Static camera with sub-pixel jitter, synthetic-game, engine mode at 3x: PSNR-Y against frame index for the baselines, the trained tiers and three 30-minute ablations}\label{fig:tsrr-static}
\end{figure}

\begin{table}[!htbp]
\caption{Static camera with sub-pixel jitter at 3x (synthetic-game, engine mode): PSNR-Y and FLIP of the last frame (index 15) of the convergence sequence. Lower FLIP is better}\label{tab:tsrr-static}
\footnotesize
\begin{tabular*}{\textwidth}{@{\extracolsep\fill}lrr}
\toprule
Method & PSNR-Y (dB) & FLIP \\
\midrule
Bilinear & \Res{tsr.static_convergence.bilinear.psnr_y_last} & \Res{tsr.static_convergence.bilinear.flip_last_frame} \\
Analytic TAA & \Res{tsr.static_convergence.analytic_taau.psnr_y_last} & \Res{tsr.static_convergence.analytic_taau.flip_last_frame} \\
\tier{igpu} & \Res{tsr.static_convergence.igpu_s3.psnr_y_last} & \Res{tsr.static_convergence.igpu_s3.flip_last_frame} \\
\tier{low} student & \Res{tsr.static_convergence.student_s3.psnr_y_last} & \Res{tsr.static_convergence.student_s3.flip_last_frame} \\
Teacher & \Res{tsr.static_convergence.teacher_s3.psnr_y_last} & \Res{tsr.static_convergence.teacher_s3.flip_last_frame} \\
Ablation reference & \Res{tsr.static_convergence.abl_reference.psnr_y_last} & \Res{tsr.static_convergence.abl_reference.flip_last_frame} \\
Ablation: no long unrolls & \Res{tsr.static_convergence.abl_no_long_unroll.psnr_y_last} & \Res{tsr.static_convergence.abl_no_long_unroll.flip_last_frame} \\
Ablation: no history & \Res{tsr.static_convergence.abl_no_history.psnr_y_last} & \Res{tsr.static_convergence.abl_no_history.flip_last_frame} \\
\botrule
\end{tabular*}
\end{table}

\FloatBarrier

\FloatBarrier

\emph{Interpretation.} On the last frame of the convergence sequence the student is \ResAbsDiff{tsr.static_convergence.student_s3.psnr_y_last}{tsr.static_convergence.analytic_taau.psnr_y_last} dB above TAA, and its FLIP error \cite{andersson2020flip} is \Res{tsr.static_convergence.student_s3.flip_last_frame} against \Res{tsr.static_convergence.analytic_taau.flip_last_frame} for TAA. \Observation{} The FLIP ordering agrees with the PSNR ordering. A single static sequence of six stills is not enough to establish a convergence rate. The decline of the model without unrolls after its peak is the motivation for the 16-frame unrolls of training, and the reference with unrolls declines much less. Both results come from one run, so the size of each difference is \Preliminary.

\subsection{Ablations}\label{sec:tsrr-ablation}

Each ablation changes one part of the reference recipe and trains for 30 minutes, the same budget as the reference. The removed components are the history (single-frame upscaling), the G-buffers (depth, motion vectors and the availability flags of the render input, so that only colour remains), the built-in motion estimator, the video samples of training and the 16-frame unrolls. \Cref{tab:tsrr-abl-abs} gives the PSNR-Y of every run on every suite and mode at 3x, and \cref{tab:tsrr-abl-delta} gives the differences to the reference. \Cref{fig:tsrr-abl} plots the differences. The ablations are short runs, so they are read against the reference and not against the main student, and differences of a few tenths of a decibel need the noise estimate given below.

\begin{table}[!htbp]
\caption{Ablations at 3x: PSNR-Y (dB) of every run on every suite and mode. Each run is a 30-minute run of the \tier{low} recipe. E: engine mode, P: plain mode. Syn: synthetic-game, Tar: TartanAir, Sin: Sintel, Vid4, Vim: Vimeo-1080p. Means over per-frame rows}\label{tab:tsrr-abl-abs}
\scriptsize
\setlength{\tabcolsep}{2pt}
\begin{tabular*}{\textwidth}{@{\extracolsep\fill}lrrrrrrrrrr}
\toprule
Run & Syn E & Syn P & Tar E & Tar P & Sin E & Sin P & Vid4 E & Vid4 P & Vim E & Vim P \\
\midrule
Reference & \Res{tsr.test_results.synthetic-game.engine.3.abl_reference.psnr_y} & \Res{tsr.test_results.synthetic-game.plain.3.abl_reference.psnr_y} & \Res{tsr.test_results.tartanair.engine.3.abl_reference.psnr_y} & \Res{tsr.test_results.tartanair.plain.3.abl_reference.psnr_y} & \Res{tsr.test_results.sintel.engine.3.abl_reference.psnr_y} & \Res{tsr.test_results.sintel.plain.3.abl_reference.psnr_y} & \Res{tsr.test_results.vid4.engine.3.abl_reference.psnr_y} & \Res{tsr.test_results.vid4.plain.3.abl_reference.psnr_y} & \Res{tsr.test_results.vimeo1080p.engine.3.abl_reference.psnr_y} & \Res{tsr.test_results.vimeo1080p.plain.3.abl_reference.psnr_y} \\
No history & \Res{tsr.test_results.synthetic-game.engine.3.abl_no_history.psnr_y} & \Res{tsr.test_results.synthetic-game.plain.3.abl_no_history.psnr_y} & \Res{tsr.test_results.tartanair.engine.3.abl_no_history.psnr_y} & \Res{tsr.test_results.tartanair.plain.3.abl_no_history.psnr_y} & \Res{tsr.test_results.sintel.engine.3.abl_no_history.psnr_y} & \Res{tsr.test_results.sintel.plain.3.abl_no_history.psnr_y} & \Res{tsr.test_results.vid4.engine.3.abl_no_history.psnr_y} & \Res{tsr.test_results.vid4.plain.3.abl_no_history.psnr_y} & \Res{tsr.test_results.vimeo1080p.engine.3.abl_no_history.psnr_y} & \Res{tsr.test_results.vimeo1080p.plain.3.abl_no_history.psnr_y} \\
No G-buffers & \Res{tsr.test_results.synthetic-game.engine.3.abl_no_gbuffers.psnr_y} & \Res{tsr.test_results.synthetic-game.plain.3.abl_no_gbuffers.psnr_y} & \Res{tsr.test_results.tartanair.engine.3.abl_no_gbuffers.psnr_y} & \Res{tsr.test_results.tartanair.plain.3.abl_no_gbuffers.psnr_y} & \Res{tsr.test_results.sintel.engine.3.abl_no_gbuffers.psnr_y} & \Res{tsr.test_results.sintel.plain.3.abl_no_gbuffers.psnr_y} & \Res{tsr.test_results.vid4.engine.3.abl_no_gbuffers.psnr_y} & \Res{tsr.test_results.vid4.plain.3.abl_no_gbuffers.psnr_y} & \Res{tsr.test_results.vimeo1080p.engine.3.abl_no_gbuffers.psnr_y} & \Res{tsr.test_results.vimeo1080p.plain.3.abl_no_gbuffers.psnr_y} \\
No motion net & \Res{tsr.test_results.synthetic-game.engine.3.abl_no_motion_net.psnr_y} & \Res{tsr.test_results.synthetic-game.plain.3.abl_no_motion_net.psnr_y} & \Res{tsr.test_results.tartanair.engine.3.abl_no_motion_net.psnr_y} & \Res{tsr.test_results.tartanair.plain.3.abl_no_motion_net.psnr_y} & \Res{tsr.test_results.sintel.engine.3.abl_no_motion_net.psnr_y} & \Res{tsr.test_results.sintel.plain.3.abl_no_motion_net.psnr_y} & \Res{tsr.test_results.vid4.engine.3.abl_no_motion_net.psnr_y} & \Res{tsr.test_results.vid4.plain.3.abl_no_motion_net.psnr_y} & \Res{tsr.test_results.vimeo1080p.engine.3.abl_no_motion_net.psnr_y} & \Res{tsr.test_results.vimeo1080p.plain.3.abl_no_motion_net.psnr_y} \\
No video training & \Res{tsr.test_results.synthetic-game.engine.3.abl_no_video_training.psnr_y} & \Res{tsr.test_results.synthetic-game.plain.3.abl_no_video_training.psnr_y} & \Res{tsr.test_results.tartanair.engine.3.abl_no_video_training.psnr_y} & \Res{tsr.test_results.tartanair.plain.3.abl_no_video_training.psnr_y} & \Res{tsr.test_results.sintel.engine.3.abl_no_video_training.psnr_y} & \Res{tsr.test_results.sintel.plain.3.abl_no_video_training.psnr_y} & \Res{tsr.test_results.vid4.engine.3.abl_no_video_training.psnr_y} & \Res{tsr.test_results.vid4.plain.3.abl_no_video_training.psnr_y} & \Res{tsr.test_results.vimeo1080p.engine.3.abl_no_video_training.psnr_y} & \Res{tsr.test_results.vimeo1080p.plain.3.abl_no_video_training.psnr_y} \\
No long unrolls & \Res{tsr.test_results.synthetic-game.engine.3.abl_no_long_unroll.psnr_y} & \Res{tsr.test_results.synthetic-game.plain.3.abl_no_long_unroll.psnr_y} & \Res{tsr.test_results.tartanair.engine.3.abl_no_long_unroll.psnr_y} & \Res{tsr.test_results.tartanair.plain.3.abl_no_long_unroll.psnr_y} & \Res{tsr.test_results.sintel.engine.3.abl_no_long_unroll.psnr_y} & \Res{tsr.test_results.sintel.plain.3.abl_no_long_unroll.psnr_y} & \Res{tsr.test_results.vid4.engine.3.abl_no_long_unroll.psnr_y} & \Res{tsr.test_results.vid4.plain.3.abl_no_long_unroll.psnr_y} & \Res{tsr.test_results.vimeo1080p.engine.3.abl_no_long_unroll.psnr_y} & \Res{tsr.test_results.vimeo1080p.plain.3.abl_no_long_unroll.psnr_y} \\
\botrule
\end{tabular*}
\end{table}
\begin{table}[!htbp]
\caption{Ablations at 3x: PSNR-Y change (dB) of each run against the reference on every suite and mode, from the same runs as the previous table. Negative values mean that removing the component costs quality. E: engine mode, P: plain mode. Syn: synthetic-game, Tar: TartanAir, Sin: Sintel, Vid4, Vim: Vimeo-1080p. }\label{tab:tsrr-abl-delta}
\scriptsize
\setlength{\tabcolsep}{2pt}
\begin{tabular*}{\textwidth}{@{\extracolsep\fill}lrrrrrrrrrr}
\toprule
Run & Syn E & Syn P & Tar E & Tar P & Sin E & Sin P & Vid4 E & Vid4 P & Vim E & Vim P \\
\midrule
No history & \ResDiff{tsr.test_results.synthetic-game.engine.3.abl_no_history.psnr_y}{tsr.test_results.synthetic-game.engine.3.abl_reference.psnr_y} & \ResDiff{tsr.test_results.synthetic-game.plain.3.abl_no_history.psnr_y}{tsr.test_results.synthetic-game.plain.3.abl_reference.psnr_y} & \ResDiff{tsr.test_results.tartanair.engine.3.abl_no_history.psnr_y}{tsr.test_results.tartanair.engine.3.abl_reference.psnr_y} & \ResDiff{tsr.test_results.tartanair.plain.3.abl_no_history.psnr_y}{tsr.test_results.tartanair.plain.3.abl_reference.psnr_y} & \ResDiff{tsr.test_results.sintel.engine.3.abl_no_history.psnr_y}{tsr.test_results.sintel.engine.3.abl_reference.psnr_y} & \ResDiff{tsr.test_results.sintel.plain.3.abl_no_history.psnr_y}{tsr.test_results.sintel.plain.3.abl_reference.psnr_y} & \ResDiff{tsr.test_results.vid4.engine.3.abl_no_history.psnr_y}{tsr.test_results.vid4.engine.3.abl_reference.psnr_y} & \ResDiff{tsr.test_results.vid4.plain.3.abl_no_history.psnr_y}{tsr.test_results.vid4.plain.3.abl_reference.psnr_y} & \ResDiff{tsr.test_results.vimeo1080p.engine.3.abl_no_history.psnr_y}{tsr.test_results.vimeo1080p.engine.3.abl_reference.psnr_y} & \ResDiff{tsr.test_results.vimeo1080p.plain.3.abl_no_history.psnr_y}{tsr.test_results.vimeo1080p.plain.3.abl_reference.psnr_y} \\
No G-buffers & \ResDiff{tsr.test_results.synthetic-game.engine.3.abl_no_gbuffers.psnr_y}{tsr.test_results.synthetic-game.engine.3.abl_reference.psnr_y} & \ResDiff{tsr.test_results.synthetic-game.plain.3.abl_no_gbuffers.psnr_y}{tsr.test_results.synthetic-game.plain.3.abl_reference.psnr_y} & \ResDiff{tsr.test_results.tartanair.engine.3.abl_no_gbuffers.psnr_y}{tsr.test_results.tartanair.engine.3.abl_reference.psnr_y} & \ResDiff{tsr.test_results.tartanair.plain.3.abl_no_gbuffers.psnr_y}{tsr.test_results.tartanair.plain.3.abl_reference.psnr_y} & \ResDiff{tsr.test_results.sintel.engine.3.abl_no_gbuffers.psnr_y}{tsr.test_results.sintel.engine.3.abl_reference.psnr_y} & \ResDiff{tsr.test_results.sintel.plain.3.abl_no_gbuffers.psnr_y}{tsr.test_results.sintel.plain.3.abl_reference.psnr_y} & \ResDiff{tsr.test_results.vid4.engine.3.abl_no_gbuffers.psnr_y}{tsr.test_results.vid4.engine.3.abl_reference.psnr_y} & \ResDiff{tsr.test_results.vid4.plain.3.abl_no_gbuffers.psnr_y}{tsr.test_results.vid4.plain.3.abl_reference.psnr_y} & \ResDiff{tsr.test_results.vimeo1080p.engine.3.abl_no_gbuffers.psnr_y}{tsr.test_results.vimeo1080p.engine.3.abl_reference.psnr_y} & \ResDiff{tsr.test_results.vimeo1080p.plain.3.abl_no_gbuffers.psnr_y}{tsr.test_results.vimeo1080p.plain.3.abl_reference.psnr_y} \\
No motion net & \ResDiff{tsr.test_results.synthetic-game.engine.3.abl_no_motion_net.psnr_y}{tsr.test_results.synthetic-game.engine.3.abl_reference.psnr_y} & \ResDiff{tsr.test_results.synthetic-game.plain.3.abl_no_motion_net.psnr_y}{tsr.test_results.synthetic-game.plain.3.abl_reference.psnr_y} & \ResDiff{tsr.test_results.tartanair.engine.3.abl_no_motion_net.psnr_y}{tsr.test_results.tartanair.engine.3.abl_reference.psnr_y} & \ResDiff{tsr.test_results.tartanair.plain.3.abl_no_motion_net.psnr_y}{tsr.test_results.tartanair.plain.3.abl_reference.psnr_y} & \ResDiff{tsr.test_results.sintel.engine.3.abl_no_motion_net.psnr_y}{tsr.test_results.sintel.engine.3.abl_reference.psnr_y} & \ResDiff{tsr.test_results.sintel.plain.3.abl_no_motion_net.psnr_y}{tsr.test_results.sintel.plain.3.abl_reference.psnr_y} & \ResDiff{tsr.test_results.vid4.engine.3.abl_no_motion_net.psnr_y}{tsr.test_results.vid4.engine.3.abl_reference.psnr_y} & \ResDiff{tsr.test_results.vid4.plain.3.abl_no_motion_net.psnr_y}{tsr.test_results.vid4.plain.3.abl_reference.psnr_y} & \ResDiff{tsr.test_results.vimeo1080p.engine.3.abl_no_motion_net.psnr_y}{tsr.test_results.vimeo1080p.engine.3.abl_reference.psnr_y} & \ResDiff{tsr.test_results.vimeo1080p.plain.3.abl_no_motion_net.psnr_y}{tsr.test_results.vimeo1080p.plain.3.abl_reference.psnr_y} \\
No video training & \ResDiff{tsr.test_results.synthetic-game.engine.3.abl_no_video_training.psnr_y}{tsr.test_results.synthetic-game.engine.3.abl_reference.psnr_y} & \ResDiff{tsr.test_results.synthetic-game.plain.3.abl_no_video_training.psnr_y}{tsr.test_results.synthetic-game.plain.3.abl_reference.psnr_y} & \ResDiff{tsr.test_results.tartanair.engine.3.abl_no_video_training.psnr_y}{tsr.test_results.tartanair.engine.3.abl_reference.psnr_y} & \ResDiff{tsr.test_results.tartanair.plain.3.abl_no_video_training.psnr_y}{tsr.test_results.tartanair.plain.3.abl_reference.psnr_y} & \ResDiff{tsr.test_results.sintel.engine.3.abl_no_video_training.psnr_y}{tsr.test_results.sintel.engine.3.abl_reference.psnr_y} & \ResDiff{tsr.test_results.sintel.plain.3.abl_no_video_training.psnr_y}{tsr.test_results.sintel.plain.3.abl_reference.psnr_y} & \ResDiff{tsr.test_results.vid4.engine.3.abl_no_video_training.psnr_y}{tsr.test_results.vid4.engine.3.abl_reference.psnr_y} & \ResDiff{tsr.test_results.vid4.plain.3.abl_no_video_training.psnr_y}{tsr.test_results.vid4.plain.3.abl_reference.psnr_y} & \ResDiff{tsr.test_results.vimeo1080p.engine.3.abl_no_video_training.psnr_y}{tsr.test_results.vimeo1080p.engine.3.abl_reference.psnr_y} & \ResDiff{tsr.test_results.vimeo1080p.plain.3.abl_no_video_training.psnr_y}{tsr.test_results.vimeo1080p.plain.3.abl_reference.psnr_y} \\
No long unrolls & \ResDiff{tsr.test_results.synthetic-game.engine.3.abl_no_long_unroll.psnr_y}{tsr.test_results.synthetic-game.engine.3.abl_reference.psnr_y} & \ResDiff{tsr.test_results.synthetic-game.plain.3.abl_no_long_unroll.psnr_y}{tsr.test_results.synthetic-game.plain.3.abl_reference.psnr_y} & \ResDiff{tsr.test_results.tartanair.engine.3.abl_no_long_unroll.psnr_y}{tsr.test_results.tartanair.engine.3.abl_reference.psnr_y} & \ResDiff{tsr.test_results.tartanair.plain.3.abl_no_long_unroll.psnr_y}{tsr.test_results.tartanair.plain.3.abl_reference.psnr_y} & \ResDiff{tsr.test_results.sintel.engine.3.abl_no_long_unroll.psnr_y}{tsr.test_results.sintel.engine.3.abl_reference.psnr_y} & \ResDiff{tsr.test_results.sintel.plain.3.abl_no_long_unroll.psnr_y}{tsr.test_results.sintel.plain.3.abl_reference.psnr_y} & \ResDiff{tsr.test_results.vid4.engine.3.abl_no_long_unroll.psnr_y}{tsr.test_results.vid4.engine.3.abl_reference.psnr_y} & \ResDiff{tsr.test_results.vid4.plain.3.abl_no_long_unroll.psnr_y}{tsr.test_results.vid4.plain.3.abl_reference.psnr_y} & \ResDiff{tsr.test_results.vimeo1080p.engine.3.abl_no_long_unroll.psnr_y}{tsr.test_results.vimeo1080p.engine.3.abl_reference.psnr_y} & \ResDiff{tsr.test_results.vimeo1080p.plain.3.abl_no_long_unroll.psnr_y}{tsr.test_results.vimeo1080p.plain.3.abl_reference.psnr_y} \\
\botrule
\end{tabular*}
\end{table}

\begin{figure}[!htbp]
\centering
\includegraphics[width=\textwidth]{\TsrOut/plots/ablations.png}
\caption{PSNR-Y change of each ablation against the 30-minute reference (dB), per suite, for engine mode (left) and plain mode (right)}\label{fig:tsrr-abl}
\end{figure}

\FloatBarrier

\emph{History.} Removing the history costs the most in engine mode, \ResDiff{tsr.test_results.synthetic-game.engine.3.abl_no_history.psnr_y}{tsr.test_results.synthetic-game.engine.3.abl_reference.psnr_y} dB on the synthetic scenes and \ResDiff{tsr.test_results.vid4.engine.3.abl_no_history.psnr_y}{tsr.test_results.vid4.engine.3.abl_reference.psnr_y} dB on Vid4. Its temporal-change error on the synthetic scenes rises from \Res{tsr.test_results.synthetic-game.engine.3.abl_reference.temporal_err} to \Res{tsr.test_results.synthetic-game.engine.3.abl_no_history.temporal_err}. \Preliminary{} In plain mode the same removal changes little on the synthetic scenes (\ResDiff{tsr.test_results.synthetic-game.plain.3.abl_no_history.psnr_y}{tsr.test_results.synthetic-game.plain.3.abl_reference.psnr_y} dB). Plain frames carry no new sub-pixel samples for the history to accumulate, so this is the expected result. \Observation{} The history helps when the input carries sub-pixel jitter, which is the mechanism of \cref{sec:tsrr-static}.

\emph{G-buffers.} In engine mode the colour-only input costs \ResDiff{tsr.test_results.synthetic-game.engine.3.abl_no_gbuffers.psnr_y}{tsr.test_results.synthetic-game.engine.3.abl_reference.psnr_y} dB on the synthetic scenes, the second largest engine effect there after the history. Plain inputs have no G-buffers, so the plain-mode differences of the same ablation come only from training. They are \ResAbsDiff{tsr.test_results.synthetic-game.plain.3.abl_no_gbuffers.psnr_y}{tsr.test_results.synthetic-game.plain.3.abl_reference.psnr_y} dB on the synthetic scenes and \ResAbsDiff{tsr.test_results.vid4.plain.3.abl_no_gbuffers.psnr_y}{tsr.test_results.vid4.plain.3.abl_reference.psnr_y} dB on Vid4, and below \ResAbsDiff{tsr.test_results.tartanair.plain.3.abl_no_gbuffers.psnr_y}{tsr.test_results.tartanair.plain.3.abl_reference.psnr_y} dB on TartanAir. \Observation{} These differences are a rough estimate of the spread between two runs of this recipe. They are smaller than the 0.3 dB run-to-run spread quoted in \cref{sec:status-legend}, but one pair of runs cannot give a reliable estimate.

\emph{Motion estimator.} In plain mode, removing the estimator costs \ResDiff{tsr.test_results.synthetic-game.plain.3.abl_no_motion_net.psnr_y}{tsr.test_results.synthetic-game.plain.3.abl_reference.psnr_y} dB on the synthetic scenes and \ResDiff{tsr.test_results.vid4.plain.3.abl_no_motion_net.psnr_y}{tsr.test_results.vid4.plain.3.abl_reference.psnr_y} dB on Vid4. Without the estimator the plain input carries no motion, so the history is warped with the zero motion vectors of the plain input, which is consistent with the loss. \Preliminary{} \UnderRevision{} The engine-mode rows of this ablation are not a clean test. Training removes the G-buffers for this run (the \texttt{ABLATIONS} entry sets \texttt{ablate\_gbuffers}), but the notebook's evaluation sets that flag only for the G-buffer ablation. The estimator run is therefore evaluated with depth and motion vectors that its training never showed, and its engine-mode differences mix the two changes. The flag must be set for this run in evaluation before its engine-mode value is used.

\emph{Video training.} Training without the video samples costs \ResDiff{tsr.test_results.synthetic-game.plain.3.abl_no_video_training.psnr_y}{tsr.test_results.synthetic-game.plain.3.abl_reference.psnr_y} dB on the synthetic scenes and \ResDiff{tsr.test_results.vid4.plain.3.abl_no_video_training.psnr_y}{tsr.test_results.vid4.plain.3.abl_reference.psnr_y} dB on Vid4 in plain mode, the largest plain-mode effect of the study. In engine mode the render-only model is \ResAbsDiff{tsr.test_results.synthetic-game.engine.3.abl_no_video_training.psnr_y}{tsr.test_results.synthetic-game.engine.3.abl_reference.psnr_y} dB higher on the synthetic scenes. \Preliminary{} The video samples teach the network how ordinary video behaves, and in this run they cost render-input quality at this budget. The article's trade-off is therefore visible here, and a larger video share or a longer budget could change it.

\emph{Long unrolls.} Removing the 16-frame unrolls costs \ResDiff{tsr.test_results.synthetic-game.plain.3.abl_no_long_unroll.psnr_y}{tsr.test_results.synthetic-game.plain.3.abl_reference.psnr_y} dB on the synthetic scenes in plain mode and \ResDiff{tsr.test_results.synthetic-game.engine.3.abl_no_long_unroll.psnr_y}{tsr.test_results.synthetic-game.engine.3.abl_reference.psnr_y} dB in engine mode. The static-scene decline of \cref{fig:tsrr-static} is the effect the unrolls target. \Preliminary{} The plain-mode loss is larger than the engine-mode loss on the synthetic scenes, and the engine-mode effect is small on every suite.

\subsection{Model tiers, cost and latency}\label{sec:tsrr-cost}

\begin{table}[!htbp]
\caption{Network tiers. Fused parameters at 3x, multiply-accumulates (GMAC) per 1080p output frame at 3x and 2x (motion estimator running), and the largest difference between the multi-branch and the fused network on random input. Trained: whether the tier was trained in this run. The parity values are reported as produced and are not established (see the text)}\label{tab:tsrr-tiers}
\footnotesize
\setlength{\tabcolsep}{3pt}
\begin{tabular*}{\textwidth}{@{\extracolsep\fill}lrrrrr}
\toprule
Tier & Trained & Fused params (3x) & GMAC 3x & GMAC 2x & Parity (3x) \\
\midrule
\tier{igpu} & Yes & \Res{tsr.model_tiers.igpu.3.fused_params} & \Res{tsr.model_tiers.igpu.3.gmac_per_frame_1080p} & \Res{tsr.model_tiers.igpu.2.gmac_per_frame_1080p} & \Res{tsr.model_tiers.igpu.3.parity_max_abs} \\
\tier{low} & Yes & \Res{tsr.model_tiers.low.3.fused_params} & \Res{tsr.model_tiers.low.3.gmac_per_frame_1080p} & \Res{tsr.model_tiers.low.2.gmac_per_frame_1080p} & \Res{tsr.model_tiers.low.3.parity_max_abs} \\
\tier{mid} & No & \Res{tsr.model_tiers.mid.3.fused_params} & \Res{tsr.model_tiers.mid.3.gmac_per_frame_1080p} & \Res{tsr.model_tiers.mid.2.gmac_per_frame_1080p} & \Res{tsr.model_tiers.mid.3.parity_max_abs} \\
\tier{high} & No & \Res{tsr.model_tiers.high.3.fused_params} & \Res{tsr.model_tiers.high.3.gmac_per_frame_1080p} & \Res{tsr.model_tiers.high.2.gmac_per_frame_1080p} & \Res{tsr.model_tiers.high.3.parity_max_abs} \\
Teacher & Yes & \Res{tsr.model_tiers.teacher.3.fused_params} & \Res{tsr.model_tiers.teacher.3.gmac_per_frame_1080p} & \Res{tsr.model_tiers.teacher.2.gmac_per_frame_1080p} & \Res{tsr.model_tiers.teacher.3.parity_max_abs} \\
\botrule
\end{tabular*}
\end{table}

\Cref{tab:tsrr-tiers} lists the tiers. The parameter counts are the trainable and the fused networks. The multiply-accumulate counts (MAC, \cref{sec:bg-cnn}) are per steady-state frame, with the motion estimator running, at 1080p output for 2x and 3x. \Confirmed{} These counts are properties of the code, so they do not vary between runs. The \tier{mid} and \tier{high} tiers were not trained in this run, so they have costs but no quality or latency. The parity column is the largest difference between the multi-branch and the fused networks on random input. \UnderRevision{} It is exactly zero for every tier, which is not expected for a fold of convolutions in FP32. The notebook's fusion function copies the network before it swaps the re-parameterised blocks, so the two sides are different graphs, but the cause of the exact zero is not established, and fusion parity is therefore not claimed here. \Cref{fig:tsrr-tiers} plots the cost of each tier at both scales.

\begin{figure}[!htbp]
\centering
\includegraphics[width=0.8\textwidth]{\TsrOut/plots/model_tiers.png}
\caption{Multiply-accumulate operations per 1080p output frame of each tier at 2x and 3x}\label{fig:tsrr-tiers}
\end{figure}

\Cref{tab:tsrr-latency} gives the measured T4 times of the trained networks, with and without engine motion vectors, next to an arithmetic estimate. The timed step includes the warps, the jitter-aware upsampling and the final blend, at batch 1, in FP16 with cuDNN autotuning and CUDA events after warm-up. \UnderRevision{} These are eager PyTorch times on a T4. They rank the networks, but they are not the latency of a deployment runtime on a target GPU, which needs the ONNX or DirectML path. The arithmetic estimate converts each MAC count to milliseconds at an assumed sustained throughput of 4 or 8 TFLOPS. It is a parameter, not a measurement, and it is not a bound. The teacher at 1080p runs in \Res{tsr.latency.teacher_s3.3.1920x1080.t4_mean_ms} ms, faster than its estimate at 4 TFLOPS of \Res{tsr.latency.teacher_s3.3.1920x1080.est_ms_assumed_igpu_4_tflops_fp16_sustained} ms. The \tier{igpu} tier runs in \Res{tsr.latency.igpu_s3.3.1920x1080.t4_mean_ms} ms, far slower than its estimate of \Res{tsr.latency.igpu_s3.3.1920x1080.est_ms_assumed_igpu_4_tflops_fp16_sustained} ms. Small networks are limited by per-operation overhead and memory traffic in eager mode, while the teacher's larger convolutions use the GPU better.

\begin{table}[!htbp]
\caption{Measured T4 time of the fused FP16 networks per output frame (eager PyTorch, batch 1, mean and 95th percentile of 100 steps) with the motion estimator running (plain video) and with engine motion vectors (estimator skipped). Mean, MV is the mean time with engine motion vectors. Est. 4 and Est. 8 are arithmetic estimates in milliseconds at an assumed sustained throughput of 4 and 8 TFLOPS. They are parameters, not measurements}\label{tab:tsrr-latency}
\footnotesize
\setlength{\tabcolsep}{2.5pt}
\begin{tabular*}{\textwidth}{@{\extracolsep\fill}llrrrrrr}
\toprule
Model & Output & Mean & P95 & Mean, MV & GMAC & Est. 4 & Est. 8 \\
\midrule
\tier{igpu} 3x & 1080p & \Res{tsr.latency.igpu_s3.3.1920x1080.t4_mean_ms} & \Res{tsr.latency.igpu_s3.3.1920x1080.t4_p95_ms} & \Res{tsr.latency.igpu_s3.3.1920x1080.t4_mean_ms_engine_mv} & \Res{tsr.latency.igpu_s3.3.1920x1080.gmac} & \Res{tsr.latency.igpu_s3.3.1920x1080.est_ms_assumed_igpu_4_tflops_fp16_sustained} & \Res{tsr.latency.igpu_s3.3.1920x1080.est_ms_assumed_low-end_dgpu_8_tflops} \\
\tier{igpu} 3x & 4K & \Res{tsr.latency.igpu_s3.3.3840x2160.t4_mean_ms} & \Res{tsr.latency.igpu_s3.3.3840x2160.t4_p95_ms} & \Res{tsr.latency.igpu_s3.3.3840x2160.t4_mean_ms_engine_mv} & \Res{tsr.latency.igpu_s3.3.3840x2160.gmac} & \Res{tsr.latency.igpu_s3.3.3840x2160.est_ms_assumed_igpu_4_tflops_fp16_sustained} & \Res{tsr.latency.igpu_s3.3.3840x2160.est_ms_assumed_low-end_dgpu_8_tflops} \\
\tier{low} 3x & 1080p & \Res{tsr.latency.student_s3.3.1920x1080.t4_mean_ms} & \Res{tsr.latency.student_s3.3.1920x1080.t4_p95_ms} & \Res{tsr.latency.student_s3.3.1920x1080.t4_mean_ms_engine_mv} & \Res{tsr.latency.student_s3.3.1920x1080.gmac} & \Res{tsr.latency.student_s3.3.1920x1080.est_ms_assumed_igpu_4_tflops_fp16_sustained} & \Res{tsr.latency.student_s3.3.1920x1080.est_ms_assumed_low-end_dgpu_8_tflops} \\
\tier{low} 3x & 4K & \Res{tsr.latency.student_s3.3.3840x2160.t4_mean_ms} & \Res{tsr.latency.student_s3.3.3840x2160.t4_p95_ms} & \Res{tsr.latency.student_s3.3.3840x2160.t4_mean_ms_engine_mv} & \Res{tsr.latency.student_s3.3.3840x2160.gmac} & \Res{tsr.latency.student_s3.3.3840x2160.est_ms_assumed_igpu_4_tflops_fp16_sustained} & \Res{tsr.latency.student_s3.3.3840x2160.est_ms_assumed_low-end_dgpu_8_tflops} \\
\tier{low} 2x & 1080p & \Res{tsr.latency.student_s2.2.1920x1080.t4_mean_ms} & \Res{tsr.latency.student_s2.2.1920x1080.t4_p95_ms} & \Res{tsr.latency.student_s2.2.1920x1080.t4_mean_ms_engine_mv} & \Res{tsr.latency.student_s2.2.1920x1080.gmac} & \Res{tsr.latency.student_s2.2.1920x1080.est_ms_assumed_igpu_4_tflops_fp16_sustained} & \Res{tsr.latency.student_s2.2.1920x1080.est_ms_assumed_low-end_dgpu_8_tflops} \\
Teacher 3x & 1080p & \Res{tsr.latency.teacher_s3.3.1920x1080.t4_mean_ms} & \Res{tsr.latency.teacher_s3.3.1920x1080.t4_p95_ms} & \Res{tsr.latency.teacher_s3.3.1920x1080.t4_mean_ms_engine_mv} & \Res{tsr.latency.teacher_s3.3.1920x1080.gmac} & \Res{tsr.latency.teacher_s3.3.1920x1080.est_ms_assumed_igpu_4_tflops_fp16_sustained} & \Res{tsr.latency.teacher_s3.3.1920x1080.est_ms_assumed_low-end_dgpu_8_tflops} \\
Teacher 3x & 4K & \Res{tsr.latency.teacher_s3.3.3840x2160.t4_mean_ms} & \Res{tsr.latency.teacher_s3.3.3840x2160.t4_p95_ms} & \Res{tsr.latency.teacher_s3.3.3840x2160.t4_mean_ms_engine_mv} & \Res{tsr.latency.teacher_s3.3.3840x2160.gmac} & \Res{tsr.latency.teacher_s3.3.3840x2160.est_ms_assumed_igpu_4_tflops_fp16_sustained} & \Res{tsr.latency.teacher_s3.3.3840x2160.est_ms_assumed_low-end_dgpu_8_tflops} \\
\botrule
\end{tabular*}
\end{table}

\begin{figure}[!htbp]
\centering
\includegraphics[width=0.8\textwidth]{\TsrOut/plots/latency.png}
\caption{Measured T4 time of the fused FP16 networks per output frame at 1080p and 4K, with the built-in motion estimator running (plain video)}\label{fig:tsrr-latency}
\end{figure}

\Cref{fig:tsrr-latency} shows the measured times. With engine motion vectors the estimator is skipped, and the time falls from \Res{tsr.latency.igpu_s3.3.1920x1080.t4_mean_ms} ms to \Res{tsr.latency.igpu_s3.3.1920x1080.t4_mean_ms_engine_mv} ms on the \tier{igpu} tier, and from \Res{tsr.latency.student_s3.3.1920x1080.t4_mean_ms} ms to \Res{tsr.latency.student_s3.3.1920x1080.t4_mean_ms_engine_mv} ms for the \tier{low} student. The 4K time of the \tier{igpu} tier is \ResRatio{tsr.latency.igpu_s3.3.3840x2160.t4_mean_ms}{tsr.latency.igpu_s3.3.1920x1080.t4_mean_ms} times its 1080p time, against four times the pixels.

\Cref{tab:tsrr-cost} compares the trained tiers and the teacher on cost and quality, on the same test sets. \Preliminary{} The \tier{low} student costs \ResRatio{tsr.model_tiers.low.3.gmac_per_frame_1080p}{tsr.model_tiers.igpu.3.gmac_per_frame_1080p} times the MAC count of the \tier{igpu} tier, and \Res{tsr.latency.student_s3.3.1920x1080.t4_mean_ms} ms against \Res{tsr.latency.igpu_s3.3.1920x1080.t4_mean_ms} ms. Its gain over the \tier{igpu} tier is \ResDiff{tsr.test_results.synthetic-game.plain.3.student_s3.psnr_y}{tsr.test_results.synthetic-game.plain.3.igpu_s3.psnr_y} dB on the synthetic scenes in plain mode and \ResDiff{tsr.test_results.vid4.plain.3.student_s3.psnr_y}{tsr.test_results.vid4.plain.3.igpu_s3.psnr_y} dB on Vid4. The teacher costs \ResRatio{tsr.model_tiers.teacher.3.gmac_per_frame_1080p}{tsr.model_tiers.low.3.gmac_per_frame_1080p} times the MAC count of the student for \ResDiff{tsr.test_results.synthetic-game.plain.3.teacher_s3.psnr_y}{tsr.test_results.synthetic-game.plain.3.student_s3.psnr_y} dB on the synthetic scenes, and for \ResDiff{tsr.test_results.vimeo1080p.plain.3.teacher_s3.psnr_y}{tsr.test_results.vimeo1080p.plain.3.student_s3.psnr_y} dB on Vimeo-1080p. \Observation{} The \tier{igpu} tier is the cheapest trained tier, and its gain over bicubic on the real plain suites is smaller than the student's, so choosing it saves cost at a visible loss of quality.

\begin{table}[!htbp]
\caption{Cost against quality at 3x. MAC per 1080p frame and measured T4 time (plain video, 1080p) against PSNR-Y (dB) on four suites and modes. The bicubic row has no network and therefore no cost}\label{tab:tsrr-cost}
\footnotesize
\setlength{\tabcolsep}{3pt}
\begin{tabular*}{\textwidth}{@{\extracolsep\fill}lrrrrrr}
\toprule
Method & GMAC & T4 ms & Synthetic, engine & Synthetic, plain & Vid4, plain & Vimeo-1080p, plain \\
\midrule
Bicubic & n/a & n/a & \Res{tsr.test_results.synthetic-game.engine.3.bicubic.psnr_y} & \Res{tsr.test_results.synthetic-game.plain.3.bicubic.psnr_y} & \Res{tsr.test_results.vid4.plain.3.bicubic.psnr_y} & \Res{tsr.test_results.vimeo1080p.plain.3.bicubic.psnr_y} \\
\tier{igpu} & \Res{tsr.model_tiers.igpu.3.gmac_per_frame_1080p} & \Res{tsr.latency.igpu_s3.3.1920x1080.t4_mean_ms} & \Res{tsr.test_results.synthetic-game.engine.3.igpu_s3.psnr_y} & \Res{tsr.test_results.synthetic-game.plain.3.igpu_s3.psnr_y} & \Res{tsr.test_results.vid4.plain.3.igpu_s3.psnr_y} & \Res{tsr.test_results.vimeo1080p.plain.3.igpu_s3.psnr_y} \\
\tier{low} student & \Res{tsr.model_tiers.low.3.gmac_per_frame_1080p} & \Res{tsr.latency.student_s3.3.1920x1080.t4_mean_ms} & \Res{tsr.test_results.synthetic-game.engine.3.student_s3.psnr_y} & \Res{tsr.test_results.synthetic-game.plain.3.student_s3.psnr_y} & \Res{tsr.test_results.vid4.plain.3.student_s3.psnr_y} & \Res{tsr.test_results.vimeo1080p.plain.3.student_s3.psnr_y} \\
Teacher & \Res{tsr.model_tiers.teacher.3.gmac_per_frame_1080p} & \Res{tsr.latency.teacher_s3.3.1920x1080.t4_mean_ms} & \Res{tsr.test_results.synthetic-game.engine.3.teacher_s3.psnr_y} & \Res{tsr.test_results.synthetic-game.plain.3.teacher_s3.psnr_y} & \Res{tsr.test_results.vid4.plain.3.teacher_s3.psnr_y} & \Res{tsr.test_results.vimeo1080p.plain.3.teacher_s3.psnr_y} \\
\botrule
\end{tabular*}
\end{table}

\emph{ONNX export.} The \tier{low} 3x student and the \tier{igpu} tier were exported as one fused recurrent step, with explicit state tensors, static shapes and opset 17, and were checked with ONNX Runtime on the CPU against PyTorch in FP32. \Cref{tab:tsrr-onnx} gives the results. \Confirmed{} The maximum difference is \Res{tsr.onnx_export.0.ort_max_abs} for the student and \Res{tsr.onnx_export.1.ort_max_abs} for the \tier{igpu} tier, so the exported graphs compute the PyTorch function on CPU to within floating-point rounding. The export was checked only on CPU in FP32, so the parity of FP16 kernels in a GPU runtime is not verified here.

\begin{table}[!htbp]
\caption{ONNX export of the fused recurrent step: file size and the largest difference between ONNX Runtime (CPU) and PyTorch (FP32) on the export input}\label{tab:tsrr-onnx}
\footnotesize
\begin{tabular*}{\textwidth}{@{\extracolsep\fill}lrr}
\toprule
Model & File size (MB) & Max difference \\
\midrule
\tier{low} 3x student & \Res{tsr.onnx_export.0.onnx_mb} & \Res{tsr.onnx_export.0.ort_max_abs} \\
\tier{igpu} 3x & \Res{tsr.onnx_export.1.onnx_mb} & \Res{tsr.onnx_export.1.ort_max_abs} \\
\botrule
\end{tabular*}
\end{table}

\subsection{Training caches and data}\label{sec:tsrr-caches}

The notebook decoded the training data once into memory-mapped caches, so that both training processes read the same files through the page cache. \Cref{tab:tsrr-caches} gives the counts. The still caches hold crops of high-resolution textures for the synthetic renderer, and the window caches hold six-frame windows with motion, validity and depth at half resolution. The caches took \Res{tsr.stage_times.caches.minutes} minutes and occupy \Res{tsr.caches.cache_gb} GB, which the notebook deletes before packaging the output. Most training windows come from Vimeo-90K septuplets \cite{xue2019vimeo}, \Res{tsr.caches.window_sources.train.vimeo_septuplet} of \Res{tsr.caches.caches.seq_train.count} windows, and the other sources are small by comparison. The notebook also prints that almost all training windows carry estimated motion from RAFT, with a small share of exact motion. That share is not written as a key, so it is not generated here.

\begin{table}[!htbp]
\caption{Training and validation caches: crops and windows, and the windows by source. Windows of the video sources carry RAFT motion. n/a: the source has no validation windows}\label{tab:tsrr-caches}
\footnotesize
\begin{tabular*}{\textwidth}{@{\extracolsep\fill}lrr}
\toprule
Item & Training & Validation \\
\midrule
Still crops (512 pixels) & \Res{tsr.caches.caches.still_train.count} & \Res{tsr.caches.caches.still_val.count} \\
Windows (6 frames at 192 pixels) & \Res{tsr.caches.caches.seq_train.count} & \Res{tsr.caches.caches.seq_val.count} \\
\midrule
Vimeo-90K septuplets windows & \Res{tsr.caches.window_sources.train.vimeo_septuplet} & \Res{tsr.caches.window_sources.val.vimeo_septuplet} \\
REDS windows & \Res{tsr.caches.window_sources.train.reds} & \Res{tsr.caches.window_sources.val.reds} \\
Vimeo-1080p windows & \Res{tsr.caches.window_sources.train.vimeo1080p} & n/a \\
GameIR windows & \Res{tsr.caches.window_sources.train.gameir} & n/a \\
TartanAir windows & \Res{tsr.caches.window_sources.train.tartanair} & n/a \\
Sintel windows & \Res{tsr.caches.window_sources.train.sintel} & n/a \\
\midrule
Cache size on disk (GB) & \multicolumn{2}{r}{\Res{tsr.caches.cache_gb}} \\
\botrule
\end{tabular*}
\end{table}

\FloatBarrier

\Cref{fig:tsrr-cache-stills} shows random still crops, the textures of the synthetic renderer, and \cref{fig:tsrr-cache-windows} shows four windows of the sequence cache. Each window panel gives the first and last frame, the motion vectors of the last frame in the notebook's colour wheel (hue is direction, brightness is magnitude) and the validity of the last frame.

\begin{figure}[!htbp]
\centering
\includegraphics[width=\textwidth]{\TsrOut/plots/cache_stills.png}
\caption{Random still crops of the training cache, the textures of the synthetic renderer}\label{fig:tsrr-cache-stills}
\end{figure}

\begin{figure}[!htbp]
\centering
\includegraphics[width=0.75\textwidth]{\TsrOut/plots/cache_windows.png}
\caption{Four windows of the sequence cache: first and last frame, motion vectors of the last frame (hue is direction, brightness is magnitude) and validity of the last frame (black is invalid)}\label{fig:tsrr-cache-windows}
\end{figure}

\Cref{fig:tsrr-cache-stills} shows textures with a wide range of detail, from photographs and paintings to rendered scenes. In \cref{fig:tsrr-cache-windows}, the motion fields of three windows are close to uniform, and the third window shows a boundary between two motions. \Observation{} The validity maps of the first two windows are black over the whole frame, which the notebook's legend reads as disoccluded or out of view, so no pixel of those windows is trusted. The temporal-change loss of training is weighted by the same validity map (\texttt{train\_tsr.py}), and the warping and temporal-change metrics are normalised by the valid area, so such a window contributes nothing to them. The third and fourth windows are valid except for small patches. The share of fully invalid windows is not in the saved outputs. Whether this is a true failure of the forward-backward check on these clips or a fault in the cache is not decided here, and it must be checked on the cache before the next run.

\subsection{Synthetic renderer preview}\label{sec:tsrr-preview}

\FloatBarrier

\Cref{fig:tsrr-preview} shows one synthetic training sequence at three frames. The jittered render input, enlarged by nearest neighbour, shows the aliasing of thin geometry, which the renderer's sub-pixel jitter is meant to let the network resolve. The plain-video input is the same frame downscaled with anti-aliasing, blur and noise, and it has no such aliasing. The target is four-sample anti-aliased. The motion-vector panel uses the colour wheel. The valid-history panel is black along the wire and fence lines and around the moving object. The depth panel separates the near and far layers, and the reactive panel shows the glow of a particle, which writes no motion vector. \Observation{} The preview confirms that the disocclusions, the thin geometry and the particles of \cref{sec:tsrr-motion} are present in the training data.

\begin{figure}[!htbp]
\centering
\includegraphics[width=\textwidth]{\TsrOut/plots/synthetic_preview.png}
\caption{One synthetic training sequence at three frames (rows): render input (jittered, nearest-neighbour enlarged), plain-video input, target, motion vectors, valid history, depth and reactive mask}\label{fig:tsrr-preview}
\end{figure}

\subsection{Visual comparisons}\label{sec:tsrr-visual}

Each visual comparison shows the last frame of each method as a crop, the history weight of the \tier{low} student, and an x-t slice, which is one image row over time. The x-t slice shows temporal flicker and ghosting as streaks. The notebook chose these clips, so they show chosen content and not a sample of the test set.

\Cref{fig:tsrr-vis-synth} shows a fast-motion synthetic sequence. The blue object on the left moves across the frame, and its edge is where the valid history is rejected. Bilinear and TAA blur the object's edge and the thin dashed wire in the lower right. The \tier{igpu} tier, the student and the teacher keep the edge sharp. The teacher draws the wire most continuously, and the student nearly as well. The history-weight panel is white over the static background and black in the crescent along the leading edge of the object, which matches the true disocclusion panel. A regular dot grid is visible in the weight map. It comes from the pixel-shuffle unfolding of the per-pixel weights (\cref{sec:tsrm}), and the output refinement removes it from the image. \Observation{} The x-t slices of all methods look alike, with no obvious extra streaks in the learned outputs, which agrees with the low temporal-change error of the student on this clip.

\Cref{fig:tsrr-vis-tar} shows a TartanAir environment in engine mode. The scene is dark, with a red sign at the top right and lit windows and lamps. Bilinear and TAA are soft throughout. The learned outputs keep the window frames and the edges of the lit sign more sharply. The history weight is dark along the strong building edges, with sparse dark dots elsewhere, and is white over most of the frame. \Observation{} The x-t slice shows the vertical structure of the sign and the lamps, with no visible trail that the history might have smeared.

\begin{figure}[!htbp]
\centering
\includegraphics[width=\textwidth]{\TsrOut/plots/visual_tartanair.png}
\caption{TartanAir held-out environment at 3x in engine mode, last-frame crops, history weight and x-t slices, in the layout of \cref{fig:tsrr-vis-synth}}\label{fig:tsrr-vis-tar}
\end{figure}

\Cref{fig:tsrr-vis-vid4} shows Vid4 as plain video. Bicubic is soft on the signs and the building edges. The learned outputs sharpen the sign lettering and the awning stripes, as in \cref{M-fig:vid4} of the article. On the cobblestones the \tier{igpu} tier, the student and the teacher add a fine texture that bicubic lacks. It is a regular pattern, not a copy of the ground truth, whose texture is irregular. The history weight shows a dot texture rather than a disocclusion structure, and the true disocclusion panel is empty for this clip. The x-t slices of the methods are close to one another, with no visible flicker. \Observation{} The cobblestone texture is the clearest difference between the methods on this clip.

The full predicted frames of each scene are shown below, one panel per method, with the ground truth last. \Cref{fig:tsrr-pred-synth} shows the synthetic frame. The largest visible differences lie along the blue object's edge and along the dashed wire in the lower right, and the background is similar across the methods. The learned outputs keep the wire more clearly than bilinear and TAA.

\begin{figure}[!htbp]
\centering
\begin{subfigure}[t]{0.32\textwidth}\centering\includegraphics[width=\linewidth]{\TsrOut/predictions/visual_synthetic_fast_bilinear.png}\caption{Bilinear}\end{subfigure}\hfill
\begin{subfigure}[t]{0.32\textwidth}\centering\includegraphics[width=\linewidth]{\TsrOut/predictions/visual_synthetic_fast_analytic_taau.png}\caption{Analytic TAA}\end{subfigure}\hfill
\begin{subfigure}[t]{0.32\textwidth}\centering\includegraphics[width=\linewidth]{\TsrOut/predictions/visual_synthetic_fast_igpu_s3.png}\caption{\tier{igpu}}\end{subfigure}\\[4pt]
\begin{subfigure}[t]{0.32\textwidth}\centering\includegraphics[width=\linewidth]{\TsrOut/predictions/visual_synthetic_fast_student_s3.png}\caption{\tier{low} student}\end{subfigure}\hfill
\begin{subfigure}[t]{0.32\textwidth}\centering\includegraphics[width=\linewidth]{\TsrOut/predictions/visual_synthetic_fast_teacher_s3.png}\caption{Teacher}\end{subfigure}\hfill
\begin{subfigure}[t]{0.32\textwidth}\centering\includegraphics[width=\linewidth]{\TsrOut/predictions/visual_synthetic_fast_ground_truth.png}\caption{Ground truth}\end{subfigure}
\caption{Full predicted frames of the synthetic fast-motion scene at 3x (last frame), one panel per method}\label{fig:tsrr-pred-synth}
\end{figure}

\FloatBarrier

\Cref{fig:tsrr-pred-tar} gives the TartanAir frames in the same layout. The learned outputs look closest to the ground truth on the red sign and the lit windows, and bilinear and TAA are visibly softer on the lit surfaces.

\begin{figure}[!htbp]
\centering
\begin{subfigure}[t]{0.32\textwidth}\centering\includegraphics[width=\linewidth]{\TsrOut/predictions/visual_tartanair_bilinear.png}\caption{Bilinear}\end{subfigure} \hfill \begin{subfigure}[t]{0.32\textwidth}\centering\includegraphics[width=\linewidth]{\TsrOut/predictions/visual_tartanair_analytic_taau.png}\caption{Analytic TAA}\end{subfigure} \hfill \begin{subfigure}[t]{0.32\textwidth}\centering\includegraphics[width=\linewidth]{\TsrOut/predictions/visual_tartanair_igpu_s3.png}\caption{\tier{igpu}}\end{subfigure}\\[4pt]
\begin{subfigure}[t]{0.32\textwidth}\centering\includegraphics[width=\linewidth]{\TsrOut/predictions/visual_tartanair_student_s3.png}\caption{\tier{low} student}\end{subfigure} \hfill \begin{subfigure}[t]{0.32\textwidth}\centering\includegraphics[width=\linewidth]{\TsrOut/predictions/visual_tartanair_teacher_s3.png}\caption{Teacher}\end{subfigure} \hfill \begin{subfigure}[t]{0.32\textwidth}\centering\includegraphics[width=\linewidth]{\TsrOut/predictions/visual_tartanair_ground_truth.png}\caption{Ground truth}\end{subfigure}
\caption{Full predicted frames of the TartanAir environment at 3x in engine mode (last frame), one panel per method}\label{fig:tsrr-pred-tar}
\end{figure}

\FloatBarrier

\Cref{fig:tsrr-pred-vid4} gives the Vid4 frames as plain video. The learned outputs sharpen the signs, and the cobblestone texture of each learned output differs from the ground truth, as described above.

\begin{figure}[!htbp]
\centering
\begin{subfigure}[t]{0.32\textwidth}\centering\includegraphics[width=\linewidth]{\TsrOut/predictions/visual_vid4_plain_bilinear.png}\caption{Bilinear}\end{subfigure} \hfill \begin{subfigure}[t]{0.32\textwidth}\centering\includegraphics[width=\linewidth]{\TsrOut/predictions/visual_vid4_plain_bicubic.png}\caption{Bicubic}\end{subfigure} \hfill \begin{subfigure}[t]{0.32\textwidth}\centering\includegraphics[width=\linewidth]{\TsrOut/predictions/visual_vid4_plain_igpu_s3.png}\caption{\tier{igpu}}\end{subfigure}\\[4pt]
\begin{subfigure}[t]{0.32\textwidth}\centering\includegraphics[width=\linewidth]{\TsrOut/predictions/visual_vid4_plain_student_s3.png}\caption{\tier{low} student}\end{subfigure} \hfill \begin{subfigure}[t]{0.32\textwidth}\centering\includegraphics[width=\linewidth]{\TsrOut/predictions/visual_vid4_plain_teacher_s3.png}\caption{Teacher}\end{subfigure} \hfill \begin{subfigure}[t]{0.32\textwidth}\centering\includegraphics[width=\linewidth]{\TsrOut/predictions/visual_vid4_plain_ground_truth.png}\caption{Ground truth}\end{subfigure}
\caption{Full predicted frames of the Vid4 calendar clip as plain video at 3x (last frame), one panel per method}\label{fig:tsrr-pred-vid4}
\end{figure}

\FloatBarrier

\subsection{Status summary and open checks}\label{sec:tsrr-summary}

The main findings of this section are listed below, each with the tag that its evidence supports.
\begin{itemize}
\item \Preliminary{} In engine mode the \tier{low} student is ahead of TAA and bicubic on every 3x suite, by the largest margin on the synthetic scenes.
\item \Preliminary{} In plain mode the \tier{low} student is ahead of bicubic on every 3x suite, by the smallest margin on Vimeo-1080p.
\item \Preliminary{} The teacher is the best 3x model on every suite and mode, at about ten times the MAC count of the student.
\item \Preliminary{} The history, the G-buffers and the plain-video training are the largest ablation effects, and the 16-frame unrolls matter for static convergence.
\item \Confirmed{} The MAC counts and parameter counts of the tiers are properties of the code.
\item \UnderRevision{} The latencies are eager PyTorch times on one T4, not the latency of a deployment runtime.
\item \UnderRevision{} The fusion parity of the tiers is exactly zero and is not established.
\item \UnderRevision{} The engine-mode rows of the motion-estimator ablation are confounded by the G-buffer flag.
\item \Observation{} Two displayed training windows have a fully invalid validity map, and the cause is open.
\end{itemize}

The results need these checks before the next run. The engine mode of the video suites uses RAFT motion from the clean frames, so it should be reported with the plain mode alongside it, or replaced. The analytic baseline should be compared with a stronger temporal baseline. The estimator ablation needs its flag set in evaluation. The validity of the training windows must be measured on the cache. The GameIR suite needs more clips, so that its LPIPS and tOF rest on more than one frame. The budgets must be long enough for the validation curves to flatten before the stage comparisons are read. The limits of the whole study and the planned revisions are in \cref{sec:status}.
